% macro conflict, first definition kept: \newcommand{\Hotz}{\mathrm{Hot}_0^{\wedge}}
\documentclass[11pt]{article}
\usepackage[no-math]{fontspec}
\providecommand\IfDocumentMetadataT[1]{}
\usepackage[english,french]{babel}
\usepackage{amsmath}
\usepackage{amssymb}
\usepackage{amsthm}
\usepackage{mathrsfs}
\usepackage{tikz-cd}
\usepackage{enumitem}

\newcommand{\printedpage}[1]{\marginpar{\footnotesize p.~#1}}
\newcommand{\msp}[1]{\marginpar{\footnotesize [#1]}}
\newcommand{\Hotz}{\mathrm{Hot}^{\wedge}_{0}}
\newcommand{\Ch}{\underline{\mathrm{Ch}}}
\newcommand{\Ty}{\underline{\mathrm{Ty}}}
\newcommand{\CNo}{C_{\mathbf{N}^{o}}}
\newcommand{\lab}[2][l]{\text{\footnotesize\begin{tabular}{@{}#1@{}}#2\end{tabular}}}
\newcommand{\Hot}{\mathrm{Hot}}
\newcommand{\Hc}{\underline{\mathrm{Hom\,cof}}}
\providecommand{\No}{}\renewcommand{\No}{\mathbf{N}^o}
\newcommand{\uHom}{\underline{\mathrm{Hom}}}
\newcommand{\comp}{\operatorname{comp}}
\newcommand{\uCh}{\underline{\mathrm{Ch}}}
\newcommand{\hot}{\mathrm{hot}}
\newcommand{\dashto}[1]{\mathrel{\overset{#1}{\rule[0.6ex]{2em}{0.4pt}}}}
\newcommand{\CCh}{\underline{\underline{\mathrm{Ch}}}}
\newcommand{\Xb}{\overline{X'}}
\newcommand{\Xbb}{\overline{\overline{X'}}}
\newcommand{\cC}{\mathcal{C}}
\newcommand{\cX}{\mathcal{X}}
\newcommand{\NN}{\mathbf{N}^{o}}
\newcommand{\Ob}{\operatorname{Ob}}
\newcommand{\CN}{\mathcal{C}_{\mathbf{N}^o}}
\input{reader_layout.tex}
\begin{document}
% ---- from Der-ChVII-456_p071-075.tex
\printedpage{69}
% head: A. Grothendieck : Dérivateurs, Ch. VII, Édition des Universités de Montpellier II et Paris VII

\msp{page 108}
Je ne trouve d'ailleurs aucune flèche directe entre $C$ et l'un des objets $C_{\infty}$, $H_{C}$; tout ce que je vois dans ce sens, c'est
\[
\begin{tikzcd}[column sep=large, row sep=huge]
\CNo \arrow[r, "{\substack{\text{foncteur de}\\ \text{localisation}}}"] \arrow[d] & C_{\infty} \\
C, &
\end{tikzcd}
\]
où la flèche verticale est la flèche de projection de $\CNo = C \times_{T} \mathbf{N}^{o}$ (avec $\mathbf{N}^{o} \xrightarrow{\ t\ } T$). Ceci n'établit une relation topologique étroite entre $C$ (en tant qu'objet de Cat\,hot ou de Hot simplement) et $C_{\infty}$, que si les deux flèches précédentes sont tout au moins des Hot-équivalences. Et le cas-type (problématique) où la question de la relation entre $C$ et $C_{\infty}$ se pose, est donc celui où $T = \mathbf{N}^{o}$, i.e. celui d'un système inductif quelconque de (petites) catégories, indexé par $\mathbf{N}$, et interprété comme défini par une catégorie fibrée $C = \CNo$ sur $\mathbf{N}^{o}$. Dans ce cas particulier, on a du moins une flèche canonique
\[
C \longrightarrow C_{\infty}
\]
sur laquelle on peut se poser la question si c'est
\msp{page 109}
une Hot-équivalence. Sinon, cela prouvera que dans le cas général, le Hot-type de $C$ (et a fortiori sa classe d'isomorphie dans Cat\,hot) n'a pas de sens géométrique intéressant. C'est à dire que $C$ est intéressante uniquement \textit{en tant que catégorie fibrée sur $T$}, permettant notamment de construire l'ind-objet $H_{T}$ de Cat\,hot ou de Hot, mais pas en tant que `catégorie nue'.

Mais je viens de me convaincre que $\CNo \longrightarrow C_{\infty}$ est bien une Hot-équivalence. Donc la flèche $\CNo \longrightarrow C$ définit une flèche \textit{dans} Hot
\[
C_{\infty} \longrightarrow C
\]
(en composant avec l'inverse dans Hot de la précédente). Est-ce un isomorphisme? Donc pour cette question, on est ramené à la question pour $\CNo \longrightarrow C$. Mais si on prend $C = T$, cela pose déjà la même question pour $\mathbf{N}^{o} \longrightarrow T$, donc la question si $T$ est asphérique, sous la seule hypothèse que $T_{\mathrm{o}}$ soit filtrante à gauche et dénombrable à l'infini (à gauche). Mais bien sûr c'est faux (p. ex. si $T$ n'a qu'un seul objet, il n'est pas pour autant asphérique !).

Cela me convainc donc qu'à défaut
% foot: Version 22 décembre 2001
\printedpage{70}
% head: A. Grothendieck : Dérivateurs, Ch. VII, Édition des Universités de Montpellier II et Paris VII
\msp{page 110}
d'hypothèse supplémentaire draconienne sur $T$, la catégorie $C$ `nue' n'offre pas d'intérêt particulier, à part son $\pi_{0}$ (qui est donné aussi par $C_{\infty}$, et par $H_{C}$). Il n'est pas dit, certes, que \underline{Ty} ne satisfasse pas à ces `conditions draconiennes'. Mais comme $C$ `nue' ne s'insère pas dans un formalisme sympathique (p.ex. on n'a pas d'interprétation géométrique de $\mathrm{Hom}_{\mathrm{Cat\,hot}}(S,C)$, ou de $\mathrm{Hom}_{\mathrm{Hot}}(S,C)$), je crois que par la suite je peux en faire abstraction. Exit a) !

\begin{flushright}
{\footnotesize 5.11.90}
\end{flushright}

J'ai quand même passé la journée entière hier à essayer de comprendre la relation de $C$ et de son type dans Hot, à celui de $C_{\infty}$. En effet, comme on a finalement une flèche dans Hot
\[
C_{\infty} \longrightarrow C,
\]
déduite du diagramme dans (Hot)
\[
\begin{tikzcd}[column sep=huge, row sep=huge]
\CNo \arrow[r, "\sim"] \arrow[d] & C_{\infty} \\
C, &
\end{tikzcd}
\]
de prouver que $C_{\infty} \longrightarrow C$ est un isomorphisme dans Hot
\msp{page 111}
semble la façon la plus simple, et peut-être unique, de prouver que l'image de $C_{\infty}$ dans Hot ne dépend pas, à isomorphisme unique près, des choix faits pour construire $C_{\infty}$. Sans compter que dans le cas motivant, $C = \Ch_{\Sigma}(X)$ ou $\Ch_{\Sigma}(X;a,b)$, $T = \Ty$, c'est la catégorie des chemins $C$ qui donne lieu à une loi de composition des chemins impeccable, fonctorielle, associative et tout, alors qu'il n'en est rien pour $C_{\infty}$. On ne trouve pas une loi de composition
\[
C_{\infty} \times C_{\infty} \longrightarrow C_{\infty}
\]
dans Cat, ni même dans Cat\,hot ou dans Hot - ou plutôt, si on arrive pourtant à construire une flèche $C_{\infty} \times C_{\infty} \longrightarrow C_{\infty}$ dans Cat, c'est encore laborieusement à coups de choix $\pm$ arbitraires, cran par cran, et sans savoir si ce qu'on obtient a une propriété d'unicité, ne serait ce que dans Hot; on ne trouve une interprétation géométrique tangible de cette composition que dans $\Hotz$ (où on a une véritable loi de groupoïde), ou dans $H_{T}$ en tant qu'objet de Ind\,Cat\,hot
% foot: Version 22 décembre 2001
\printedpage{71}
% head: A. Grothendieck : Dérivateurs, Ch. VII, Édition des Universités de Montpellier II et Paris VII
\msp{page 112}
ou de Ind\,Hot. (Donc la limite inductive dans $\Hotz$ n'est autre que l'élément de $\Hotz$ défini par $C_{\infty}$.)

Mes efforts pour prouver que
\[
C_{\infty} \longrightarrow C
\]
était un isomorphisme dans Hot, ou ne serait-ce qu'un 1-isomorphisme (i.e. un isomorphisme non seulement pour les $\pi_{0}$, ce qui est toujours vrai, mais aussi pour les $\pi_{1}$ de toutes les composantes), ont échoués. J'ai trouvé qu'il suffisait que le foncteur % [sic] ont échoués -> ont échoué, accord fautif du participe
\[
T \longrightarrow T_{\mathrm{o}}
\]
soit ou bien lisse et coasphérique, ou bien propre et asphérique - donc tout est OK si $T = T_{\mathrm{o}}$, i.e. si $T$ est ordonnée. Mais dans le cas $T = \Ty$, le foncteur $T \longrightarrow T_{\mathrm{o}}$ n'est ni propre, ni lisse (pas même 0-propre ou 0-lisse).

Finalement j'en suis arrivé à la conviction que dans le cas qui nous occupe ($T = \Ty$, $C = \Ch_{\Sigma}(X;a,b)$), la flèche $C_{\infty} \longrightarrow C$ n'est en général \textit{pas} un isomorphisme dans Hot. En effet, dans le cas $a = b$, si cela \texttt{[était vrai]},
\msp{page 113}
on aurait une description du Hot-groupe $\Omega(X,a) = \Ch(X,a)$ (\footnote[23]{On fait ici $\Sigma = \mathrm{Fl}(X)$.}) en termes d'une catégorie $C$ munie d'une loi de composition \textit{associative} (et unitaire) \textit{dans} Cat, redonnant la loi de groupe de $\Omega$ en passant de (Cat) à Hot. Prenant le nerf de $C$, on trouverait un monoïde semi-simplicial qui, dans Hot, donnerait la loi de groupe de $\Omega(X,a)$. Or je me rappelle (\footnote[24]{C'est un erreur - ça doit être pour d'\textit{autres} objets groupe de Hot que les groupes de lacets. J'ai vu dans Quillen que $\Omega(X,x)$ doit se décrire par un groupe semi-simplicial (quand $X$ est un ensemble semi-simplicial, $x$ un sommet). Donc le problème si $C_{\infty} \simeq C$ dans Hot reste entier.}) qu'on m'a expliqué qu'en général, la loi de groupe de $\Omega(X,a)$ ne pourrait pas être obtenue par une loi \textit{associative} sur un ensemble semi-simplicial. Ainsi, l'insuccès de mes efforts semble tenir, non à ma maladresse, mais à la nature des choses. % [sic] C'est un erreur (note 24) -> C'est une erreur, genre

Je vais résumer les faits positifs auxquels je suis arrivé dans le diagramme suivant, qui a lieu dans la situation générale $C \longrightarrow T$ envisagée au début du présent paragraphe :

\msp{page 114}
% foot: Version 22 décembre 2001
\printedpage{72}
% head: A. Grothendieck : Dérivateurs, Ch. VII, Édition des Universités de Montpellier II et Paris VII
\[
\begin{tikzcd}[column sep=small, row sep=huge]
 & & C_{t'} \arrow[d, "{u^{*}\quad\lab{(pour $u : t \longrightarrow t'$\\ dans $T$)}}"] \arrow[ldd, dotted] \arrow[rdd, dotted] & \\
 & & C_{t} \arrow[ld, swap, "{\lab{flèche définie dans\\ Ind\,Cat\,hot\\ (\texttt{[ou dans]}\\ IndHot,\\ Ind\,($\Hotz$))}}"] \arrow[rd, "{\lab[r]{flèche dans Cat\\ (donc dans\\ Cat\,hot,\\ Hot,\\ $\Hotz$)}}"] & \\
\varinjlim C_{t} = & |[draw, circle]| H_{T} \arrow[rd, swap, "{\lab{flèche\\ définie dans\\ Ind\,Cat\,hot\\ ou dans Ind\,Hot\\ ou dans Ind\,$\Hotz$}}"] & \lab{diagramme\\ commutatif\\ dans Ind\,Hot\\ ou Ind\,($\Hotz$)} & C \\
\lab{dans Ind\,Cat\,hot\\ (ou dans Ind\,Hot\\ ou dans Ind\,($\Hotz$))} & & |[draw, circle]| C_{\infty} \arrow[ru, swap, "{\lab[r]{flèche\\ définie\\ dans Hot,\\ donc dans $\Hotz$,\\ en général pas un\\ isomorphisme, même\\ dans $\Hotz$}}"] & \\
 & & = \varinjlim H_{T} \text{ dans } \Hotz &
\end{tikzcd}
\]

{\small\texttt{[Les annotations dans ce diagramme concernent les flèches en traits pleins.]}}

Les objets importants sont $H_{T}$ et $C_{\infty}$, alors que $C$ (sur $T$) et les $C_{t}$ prennent un rôle technique intermédiaire, pour construire $H_{T}$ et $C_{\infty}$. On peut considérer $C_{\infty}$ comme une catégorie (objet de (Cat)), mais dont la construction dépend de nombreux arbitraires, à l'inverse de $H_{T}$ dans Ind\,Cat\,hot. C'est l'objet de $\Hotz$ déduit de $C_{\infty}$ qui a un sens canonique, et peut être décrit comme la limite inductive de $H_{T}$ comme objet de Ind\,$\Hotz$, i.e. comme représentant le même foncteur sur $\mathrm{Hot}_{0}$ que $H_{T}$.

\msp{page 115}
Dans le cas géométrique $T = \Ty$, et $C$ de la forme $\Ch_{\Sigma}(X;a,b)$, donnant lieu à $C_{\infty}(a,b)$, on \textit{n'a pas} de définition naturelle d'un bifoncteur de composition des chemins
\[
C_{\infty}(a,b) \times C_{\infty}(b,c) \ \longrightarrow\ C_{\infty}(a,c).
\]
On peut définir, certes, des accouplements
\[
\Ch_{\Sigma,\tau}(a,b) \times \Ch_{\Sigma,\tau'}(b,c) \ \longrightarrow\ \Ch_{\Sigma,\tau\tau'}(a,c),
\]
ou, avec les notations relatives à la donnée de
\[
\begin{aligned}
&\underline{t} : \mathbf{N}^{o} \longrightarrow \Ty,\\
&\Ch_{n}(a,b) \times \Ch_{m}(b,c) \longrightarrow \Ch_{n+m}(a,c),
\end{aligned}
\]
% foot: Version 22 décembre 2001
\printedpage{73}
% head: A. Grothendieck : Dérivateurs, Ch. VII, Édition des Universités de Montpellier II et Paris VII
mais ces accouplements \textit{ne sont pas} compatibles avec les morphismes de transition des trois systèmes inductifs de catégories $(\Ch_{n})_{n\in\mathbf{N}}$, ils ne le sont \textit{qu'à homotopie près.} Cela suffit cependant pour avoir des accouplements parfaits
\[
H_{T}(a,b) \times H_{T}(b,c) \ \longrightarrow\ H_{T}(a,c),
\]
avec associativité, et (si $\Sigma = \mathrm{Fl}(X)$ tout entier) existence de \textit{vrais} inverses, i.e. de flèches (dans Ind\,Cat\,hot) de `passage au chemin opposé'

\msp{page 116}
\[
H_{T}(a,b) \ \longrightarrow\ H_{T}(b,a), \qquad c \longmapsto c^{o}, \tag*{$(*)$}
\]
donnant lieu au diagramme commutatif exprimant que c'est bien un inverse pour la composition. Une autre façon de le dire, c'est que pour tout objet $Y$ de Cat\,hot (ou de Ind\,Cat\,hot), les ensembles
\[
\mathrm{Hom}_{\mathrm{Cat\,hot}}(Y, H_{T}(a,b)) \overset{\text{déf}}{=} H_{T}^{Y}(a,b)
\]
forment, pour $a,b \in X$ variables, les flèches d'un \textit{groupoïde} (\footnote[25]{Si $\Sigma = \mathrm{Fl}(X)$. Si $\Sigma$ est plus petit, au lieu d'un groupoïde, on trouve une \textit{catégorie} sans plus.}) $\mathrm{Ch}_{\Sigma}^{Y}(X)$ ayant pour ensemble de sommets $\mathrm{Ob}\,X$, et $H_{T}^{Y}(a,b)$ étant formé des flèches de $a$ vers $b$. L'application
\[
H_{T}^{Y}(a,b) \longrightarrow H_{T}^{Y}(b,a), \qquad c \longmapsto c^{o},
\]
déduite du foncteur $(*)$ plus haut n'est autre que l'application de passage à l'inverse. Cette catégorie n'est d'ailleurs autre que celle déduit de % [sic] celle déduit -> celle déduite, accord
\[
X(Y)\Sigma_{Y}^{-1}, \qquad \text{où } \left\{
\begin{aligned}
&X(Y) = \underline{\mathrm{Hom}}_{\mathrm{Cat}}(Y,X)\\
&\Sigma_{Y} \subseteq \mathrm{Fl}(X(Y)),\\
&\text{défini en termes de } \Sigma \text{ `argument par argument'},
\end{aligned}
\right.
\]
à l'aide du foncteur
\[
X \longrightarrow X(Y), \qquad a \longmapsto a_{Y},
\]
\msp{page 117}
associant à tout $a \in X$ le foncteur \textit{constant} sur $Y$ de valeur $a$, en prenant
\[
H_{T}(a,b) \simeq \mathrm{Hom}_{X(Y)\Sigma_{Y}^{-1}}(a_{Y}, b_{Y}) \ .
\]

Dans le même cas géométrique, je me suis convaincu d'ailleurs que $C_{\infty}$ était, en tant qu'objet de Hot (pas seulement de $\Hotz$) \textit{unique à} Hot-\textit{isomorphisme près,} mais (chose étrange) peut-être pas à Hot-isomorphisme \textit{unique} près. De façon générale, revenant aux
% foot: Version 22 décembre 2001

% ---- from Der-ChVII-456_p076-080.tex
\printedpage{74}
% head: A. Grothendieck : Dérivateurs, Ch. VII, Édition des Universités de Montpellier II et Paris VII
conditions du début du paragraphe, $C_\infty$ doit s'écrire $C_\infty(\underline{t})$, elle dépend du choix du foncteur cofinal
\[
\underline{t} : \No \longrightarrow T.
\]
Quand $\underline{t}$ varie dans la catégorie
\[
\Hc(\No, T),
\]
les $C_\infty(\underline{t})$ varient de façon contravariante en $\underline{t}$ (donc définissent une catégorie fibrée sur $\Hc(\No, T)$). De plus, le foncteur
\[
C_\infty(\underline{t'}) \longrightarrow C_\infty(\underline{t})
\]
associé à un morphisme de systèmes projectifs
\[
\underline{t} \longrightarrow \underline{t'}
\]
est un \textit{isomorphisme} dans Hot. Donc le
\msp{page 118}
foncteur
\[
\underline{t} \longmapsto C_\infty(\underline{t}), \qquad \Hc(\No, T) \longrightarrow (\Hot)
\]
se factorise par le groupoïde fondamental
\[
\Pi(\Hc(\No, T)) \longrightarrow (\Hot).
\]
Quand ce groupoïde, i.e. quand la catégorie $\Hc$ elle-même est connexe, on peut donc conclure que tous les $C_\infty(\underline{t})$ sont isomorphes entre eux dans Hot. Pour conclure même à un système \textit{transitif} d'isomorphismes entre les $C_\infty(\underline{t})$, i.e. que $C_\infty(\underline{t})$ est ‘unique à isomorphisme \textit{unique} près’, il faudrait savoir de plus que le groupoïde fondamental est équivalent à la catégorie ponctuelle (groupe fondamental trivial), ou, comme on dit, que $\Hc$ soit \textit{simplement connexe,} ou 1-connexe.

La connexité de $\Hc$ est assurée quand la catégorie $T$ a (comme c'est le cas pour $\underline{\mathrm{Ty}}$) une loi de composition $(t,t') \longmapsto t * t'$, avec des morphismes fonctoriels en $t, t'$
\[
\begin{tikzcd}[row sep=small]
 & t \\
t * t' \arrow[ru] \arrow[rd] & \\
 & t'
\end{tikzcd}
\]
(cette dernière condition étant automatique
% foot: Version 22 décembre 2001
\printedpage{75}
% head: A. Grothendieck : Dérivateurs, Ch. VII, Édition des Universités de Montpellier II et Paris VII
\msp{page 119}
si la loi de composition admet l'objet final de $T$, supposé existant, comme unité bilatère). Dans ce cas, $\underline{t} = (t_n)_n$ et $\underline{t'} = (t'_n)_n$ sont reliés entre eux par $\underline{t} * \underline{t'} = (t_n * t'_n)_n$, puisqu'on a un morphisme de systèmes projectifs
\[
\begin{tikzcd}[column sep=small]
 & \underline{t} * \underline{t'} \arrow[ld] \arrow[rd] & \\
\underline{t} & & \underline{t'}.
\end{tikzcd}
\]
Je ne sais d'ailleurs pas si la flèche dans Hot
\[
C_\infty(\underline{t}) \longrightarrow C_\infty(\underline{t'})
\]
associée à $\underline{t'} \longrightarrow \underline{t}$ dans $\Hc$ ne dépend que de $\underline{t}$, $\underline{t'}$, et pas du choix de $\underline{t'} \longrightarrow \underline{t}$. Ce qui est clair, c'est que son image dans $\Hotz$ est unique, c'est l'identité quand on identifie $C_\infty(\underline{t})$ et $C_\infty(\underline{t'})$ au seul et même élément de $\Hotz$ défini par $\varinjlim H_T$.

Et pourtant, je m'aperçois à l'instant que, sous l'hypothèse de l'existence de la loi de composition et de son unité bilatère, et admettant de plus que la loi en question est associative à isomorphisme près, \textit{on a transitivité} dans les isomorphismes (canoniques!) entre les $C_\infty(\underline{t})$ qu'on vient d'indiquer !

\msp{page 120}
(Sans avoir pour autant à calculer le groupe fondamental de $\Hc(\No, T)$!) Si on a trois éléments $\underline{t}$, $\underline{t'}$, $\underline{t''}$ de $\Hc$, il suffit de regarder le diagramme déduit du diagramme commutatif suivant
\[
\begin{tikzcd}[column sep=small]
 & & \underline{t} * \underline{t'} * \underline{t''} \arrow[ld] \arrow[rd] & & \\
 & \underline{t} * \underline{t'} \arrow[ld] \arrow[rd] & & \underline{t'} * \underline{t''} \arrow[ld] \arrow[rd] & \\
\underline{t} & & \underline{t'} & & \underline{t''}
\end{tikzcd}
\]
en lui appliquant le contrafoncteur $\underline{s} \longrightarrow C_\infty(\underline{s})$ de $\Hc$ dans Hot.

Donc dans le cas des $C(a,b) = \underline{\mathrm{Ch}}_\Sigma(X; a, b)$, je n'ai plus aucun scrupule à considérer la catégorie $C_\infty(a,b)$ comme un objet \textit{canonique} de Hot, relié à l'ind-objet canonique $H_T(a,b)$ dans Hot par un homomorphisme canonique
\[
C_\infty(a,b) \longleftarrow H_T(a,b), \tag*{$(*)$}
\]
qui dans $\Hotz$ donne un isomorphisme
\[
C_\infty(a,b) \simeq \varinjlim H_T(a,b), \tag*{$(**)$}
\]
i.e. la flèche $(*)$ induit, pour $\xi \in \mathrm{Ob}\,\Hot$, une application ensembliste
\[
\mathrm{Hom}_{\Hot}(\xi, C_\infty(a,b)) \longleftarrow \underbrace{\mathrm{Hom}_{\mathrm{Ind}\,\Hot}(\xi, H_T(a,b))}_{\underset{n}{\varinjlim}\, \mathrm{Hom}_{\Hot}(\xi, C_n(a,b))},
\]
qui est \textit{bijective} si $\xi$ est dans Hot.
% foot: Version 22 décembre 2001

\printedpage{76}
% head: A. Grothendieck : Dérivateurs, Ch. VII, Édition des Universités de Montpellier II et Paris VII
\msp{page 121}
Bien entendu, les flèches
\[
C_\infty(\underline{t'}) \longrightarrow C_\infty(\underline{t}) \tag*{$(*)$}
\]
dans Hot, associés aux flèches $\underline{t} \longrightarrow \underline{t'}$ dans $\Hc(\No, T)$, sont en fait des flèches dans % [sic] associés -> associées, accord avec « les flèches »
(Cat) elle-même - et même des foncteurs pleinement fidèles, dans le cas ‘géométrique’ et pour les deux flèches canoniques
\[
\begin{tikzcd}[column sep=small]
 & \underline{t} * \underline{t'} \arrow[ld] \arrow[rd] & \\
\underline{t} & & \underline{t'}.
\end{tikzcd}
\]
Elles sont isomorphiques \textit{dans} Hot, parce qu'elles définissent un homomorphisme sur les ind-objets associés de Cat\,hot, donc aussi pour les ind-objets associés dans Hot et dans $\Hotz$, d'où résulte que
\[
\mathrm{Hom}_{\Hot}(Y, C_\infty(\underline{t'})) \longrightarrow \mathrm{Hom}_{\Hot}(Y, C_\infty(\underline{t}))
\]
associé à $(*)$ est une bijection si $Y$ \textit{est finie} (\footnote[26]{Ce qui suffit à assurer que $(*)$ est \texttt{[un]} isomorphisme dans Hot : $\Hot \longrightarrow \Hotz$ est conservatif!}), cette flèche s'identifie alors à
\[
\underset{n}{\varinjlim}\, \mathrm{Hom}_{\Hot}(Y, C_{t'_n}) \longrightarrow \underset{n}{\varinjlim}\, \mathrm{Hom}_{\Hot}(Y, C_{t_n}).
\]
(Car $Z \longmapsto \mathrm{Hom}_{\Hot}(Y, Z)$ commute aux limites inductives filtrantes, si $Y$ est une catégorie \textit{finie.} Cela provient du fait qu'une $\varinjlim$ filtrante de catégories de Thomason-Kan est itou \dots) Mais j'ignore par contre si $(*)$ est \texttt{[un]} isomorphisme dans Cat\,hot. Seulement que pour une catégorie \textit{finie} $Y$, elle induit
\[
\mathrm{Hom}_{\mathrm{Cat\,hot}}(Y, C_\infty(\underline{t'})) \xrightarrow{\ \sim\ } \mathrm{Hom}_{\mathrm{Cat\,hot}}(Y, C_\infty(\underline{t})),
\]
les deux membres s'identifiant alors à
\msp{page 122}
\[
H_T(Y)_{\mathrm{Cat\,hot}} = \mathrm{Hom}_{\mathrm{Ind\,Cat\,hot}}(Y, H_T).
\]
Pour résumer ces lumières et perplexités sur les $C_\infty(\underline{t})$ :

\begin{enumerate}
\item[a)] À tout $\underline{t} = (t_n)_{n \in \No} \in \Hc(\No, T)$ est associé un système inductif, indexé par $\mathbf{N}$, de catégories $C_n(\underline{t})$ $(= C_{t_n})$, et leur limite inductive
\[
C_\infty(\underline{t}).
\]
% foot: Version 22 décembre 2001

\printedpage{77}
% head: A. Grothendieck : Dérivateurs, Ch. VII, Édition des Universités de Montpellier II et Paris VII
\item[b)] Ce système inductif, et aussi sa limite est \texttt{[un]} foncteur contravariant en $\underline{t}$. Si $\underline{t}, \underline{t'} \in \mathrm{Ob}\,\Hc$, alors on peut comparer $C_\infty(\underline{t})$, $C_\infty(\underline{t'})$ et les systèmes inductifs correspondants, grâce à
\[
\begin{tikzcd}[column sep=small]
 & \underline{t} * \underline{t'} \arrow[ld] \arrow[rd] & \\
\underline{t} & & \underline{t'}.
\end{tikzcd}
\]

\item[c)] L'objet
\[
(C_n(\underline{t}))_{n \in \mathbf{N}} \in \mathrm{Ob\,Ind\,Cat\,hot}
\]
défini par le système inductif $(C_n(\underline{t}))_{n \in \mathbf{N}}$ de Cat est (à isomorphisme unique près) indépendant du choix de $\underline{t}$, soit $H_C$ ou $H$. Les homomorphismes
\[
C_\infty(\underline{t'}) \xrightarrow{u^*} C_\infty(\underline{t}) \qquad \text{associé à } u : \underline{t} \longrightarrow \underline{t'} \tag*{$(*)$}
\]
sont compatibles avec cette donnée, i.e. on a commutativité dans
\[
\begin{tikzcd}
\underset{\mathrm{Ind\,Cat\,hot}}{\varinjlim} C_*(\underline{t'}) \arrow[rr] & & \underset{\mathrm{Ind\,Cat\,hot}}{\varinjlim} C_*(\underline{t}) \\
 & H \arrow[lu, "\sim"] \arrow[ru, "\sim"'] &
\end{tikzcd}
\]

\msp{page 123}
\item[d)] Il s'ensuit que les homomorphismes $(*)$ sont des \textit{isomorphismes} dans Hot, donc tous les $C_\infty(\underline{t})$ sont isomorphes \textit{dans} Hot. L'astuce des $t * t'$ (dans le cas où on a une loi $t * t'$ associative et admettant l'objet final comme objet unité) donne même des isomorphismes \textit{canoniques}
\[
C_\infty(\underline{t'}) \longrightarrow C_\infty(\underline{t}) \qquad \text{dans Hot,}
\]
et c'est là un système transitif d'isomorphismes. Donc on a dans Hot un objet canonique
\[
C_\infty \in \mathrm{Ob}\,\Hot,
\]
dont la relation à $H \in \mathrm{Ob\,Ind\,Hot}$ est déduite de c), c'est une flèche canonique
\[
H \longrightarrow C_\infty \qquad \text{dans Ind\,Hot (\footnotemark[27]).}
\]
\footnotetext[27]{\textbf{NB} $H$ n'est pratiquement jamais essentiellement constant, donc la flèche ci-contre \texttt{[n'est]} pas un isomorphisme dans Ind\,Hot.}

Cette flèche induit un isomorphisme pour les objets correspondants de $\Hotz$; mais en regardant $H$ non comme ind-objet de $\Hotz$, mais prenant le foncteur $H_0$ induit sur $\Hot_0$ par le foncteur ind-représentable $H$ sur Hot. Ce foncteur induit est ‘représentable’ par l'objet $C_\infty$ de la catégorie plus grande
\[
H_0 \simeq C_\infty \qquad \text{dans } \Hotz.
\]

(Si le foncteur canonique $\Hot \longrightarrow \Hotz$ est pleinement fidèle, alors cette formule décrit $C_\infty$ comme objet de Hot, à isomorphisme unique près, en termes de $H_0$, donc aussi en termes de $H \in \mathrm{Ob\,Ind\,Cat\,hot}$.)
% foot: Version 22 décembre 2001

\printedpage{78}
% head: A. Grothendieck : Dérivateurs, Ch. VII, Édition des Universités de Montpellier II et Paris VII
\msp{page 124}
\item[e)] J'ignore par contre si les flèches
\[
C_\infty(\underline{t'}) \longrightarrow C_\infty(\underline{t})
\]
sont des isomorphismes dans Cat\,hot, bien qu'elles soient définies en termes d'homomorphismes de systèmes inductifs de catégories, qui sont des isomorphismes dans Ind\,Cat\,hot. Il en est ainsi, même pour les flèches $C_\infty(\underline{t}), C_\infty(\underline{t'}) \rightrightarrows C_\infty(\underline{t} * \underline{t'})$ utilisées pour comparer $C_\infty(\underline{t})$ et $C_\infty(\underline{t'})$ canoniquement. Ainsi, je ne vois pas comment définir des flèches en sens inverse dans Cat, qui seraient candidats à être des isomorphismes inverses dans Cat\,hot.

\item[f)] Dans le cas géométrique ($T = \underline{\mathrm{Ty}}$, situation décrite en termes de $X$, $\Sigma$, et $a, b \in \mathrm{Ob}\,X$, d'où des $C_\infty(\underline{t}; a, b)$, $C_\infty(\overline{a, b})$, $H(a,b)$ etc.), on \textit{n'a pas} d'opération de composition de chemins au niveau des \textit{catégories} $C_\infty(\underline{t}; a, b)$ et $C_\infty(\underline{t}; b, c)$, mais seulement au niveau des objets $H(a,b)$, $H(b,c)$ de Ind\,Cat\,hot ou de Ind\,Hot. Suis-je seulement sûr de savoir définir \textit{dans} Hot
\[
C_\infty(a,b) \times C_\infty(b,c) \longrightarrow C_\infty(a,c) \ \ ?
\]

Sûrement, il n'y a pas de problèmes pour les images dans $\Hotz$, puisque là ce sont les mêmes que celles de $H(a,b)$, $H(b,c)$, $H(a,c)$ respectivement, et on a des accouplements parfaits, associatifs, avec unité, et avec inverses là ou il doit y % [sic] là ou -> là où, accent omis sur le relatif
\msp{page 125}
en avoir. Mais pas pour les $C_\infty(a,b)$ eux-mêmes?
\end{enumerate}

C'est un peu vexant qu'avec tout ce travail il semblerait qu'on n'ait pas même été foutu de définir l'analogue dans Cat du classique espace des lacets $\Omega(X, x)$ comme un groupe \textit{dans} Hot. (Que fait \textsc{Thomason}?) Mais voilà une idée pensable [?]. Pour définir \textit{quand-même}
\[
C_\infty(a,b) \times C_\infty(b,c) \longrightarrow C_\infty(a,c)
\]
dans Hot, on va décrire les deux foncteurs $C_\infty(a,b)$ et $C_\infty(b,c)$ par des systèmes projectifs différents $\underline{t}$, $\underline{t'}$. On peut l'expliciter en partant d'un $\tau_0$ de longueur utile $\geq 2$, et en prenant dans les deux cas
\[
t_n = t'_n = \tau_0^n,
\]
mais des morphismes de transition différents
\[
\begin{aligned}
(t_{n+p} \longrightarrow t_n) &= \pi_p * \mathrm{id}_{\tau_0^n} \\
(t'_{n+p} \longrightarrow t'_n) &= \mathrm{id}_{\tau_0^n} * \pi_p,
\end{aligned}
\]
où $\pi_p : \tau_0^p \longrightarrow 1$ est la dégénérescence totale. On définira $C_\infty(a,c)$ en termes de $\underline{t''} = \underline{t} * \underline{t'}$, donc
\[
\left\{
\begin{aligned}
t''_n &= \tau_0^{2n} \\
(t''_{n+p} \longrightarrow t''_n) &= \pi_p * \mathrm{id}_{\tau_0^{2n}} * \pi_p.
\end{aligned}
\right.
\]
% foot: Version 22 décembre 2001

% ---- from Der-ChVII-456_p081-085.tex
\printedpage{79}
% head: A. Grothendieck : Dérivateurs, Ch. VII, Édition des Universités de Montpellier II et Paris VII
% foot: Version 22 décembre 2001

Cette fois, les foncteurs de composition
\[ C_n(\underline{t};a,b)\times C_n(\underline{t'};b,c) \longrightarrow C_n(\underline{t''};a,c), \]
sont compatibles avec les foncteurs transition pour $n$ variable, et on trouve
\msp{page 126}
\[ C_\infty(\underline{t};a,b)\times C_\infty(\underline{t'};b,c) \longrightarrow C_\infty(\underline{t''};a,c), \]
donc une flèche dans Hot
\[ \comp_{\tau_0} : C_\infty\times C_\infty \longrightarrow C_\infty. \]
Il faut voir qu'elle ne dépend pas du choix de $\tau_0$, mais si on a $\tau_0$ et $\tau'_0$, on les compare par $\tau_0\times\tau'_0 \begin{array}{l}\longrightarrow \tau_0\\ \longrightarrow \tau'_0\end{array}$, il suffit donc de voir que s'il existe une flèche $\tau_0\longrightarrow\tau'_0$, alors $\comp_{\tau_0} = \comp_{\tau'_0}$. Mais une telle flèche induit des flèches $\underline{t}_{\tau_0}\longrightarrow\underline{t}_{\tau'_0}$, $\underline{t'}_{\tau_0}\longrightarrow\underline{t'}_{\tau'_0}$, $\underline{t''}_{\tau_0}\longrightarrow\underline{t''}_{\tau'_0}$, d'où un diagramme dans Cat
\[
\begin{tikzcd}
C_\infty(\underline{t}_{\tau'_0};a,b)\times C_\infty(\underline{t'}_{\tau'_0};b,c) \arrow[r, "{\comp_{\tau'_0}}"] \arrow[d] & C_\infty(\underline{t''}_{\tau'_0};a,c) \arrow[d] \\
C_\infty(\underline{t}_{\tau_0};a,b)\times C_\infty(\underline{t'}_{\tau_0};b,c) \arrow[r, "{\comp_{\tau_0}}"] & C_\infty(\underline{t''}_{\tau_0};a,c),
\end{tikzcd}
\]
dont on vérifie tout de suite qu'il est commutatif. Or ce diagramme peut se réécrire dans Hot comme
\[
\begin{tikzcd}
C_\infty(a,b)\times C_\infty(b,c) \arrow[r, "{\comp_{\tau'_0}}"] \arrow[d, "{\mathrm{id}}"] & C_\infty(a,c) \arrow[d, "{\mathrm{id}}"] \\
C_\infty(a,b)\times C_\infty(b,c) \arrow[r, "{\comp_{\tau_0}}"] & C_\infty(a,c),
\end{tikzcd}
\]
\msp{page 127}
d'où $\comp_{\tau_0} = \comp_{\tau'_0}$.

Je recule devant la vraie difficulté - celle de prouver l'associativité de cette loi de composition. Elle devient associative dans $\Hot^{\wedge}_{0}$ - mais pour que deux flèches dans Hot soient égales (ici donc
\[ C_\infty(a,b)\times C_\infty(b,c)\times C_\infty(c,d) \rightrightarrows C_\infty(a,d)\quad ), \]
suffit-il qu'il en soit ainsi de leurs images dans $\Hot^{\wedge}_{0}$ - en d'autres termes, le foncteur
\[ \Hot\longrightarrow\Hot^{\wedge}_{0}, \]
est-il (en plus d'être conservatif) au moins fidèle (sinon pleinement fidèle)?

\clearpage
\printedpage{80}
% head: A. Grothendieck : Dérivateurs, Ch. VII, Édition des Universités de Montpellier II et Paris VII
% foot: Version 22 décembre 2001

Il y a aussi que cette loi de composition est compatible avec celle des $H(a,b)$ dans Ind\,Hot, pour les flèches canoniques
\[ H(a,b)\longrightarrow C_\infty(a,b) \qquad \text{dans Ind\,Hot}, \]
si on saurait que ce sont des épimorphismes, et le restent par produits cartésiens, on aurait gagné.

Bien sûr, on saura que c'est associatif, une fois qu'on saura que la composititon définie ici est compatible avec la composition des chemins dans le cadre semi-simplicial. Mais comment fait-on dans ce cadre, pour composer les chemins de longueur non bornée? N'y a-t-il pas là exactement la même difficulté? À moins de s'en % [sic] composititon -> composition, misspelling % Reading resolved against the supplied PDF; see corrections.csv.
\msp{page 128}
tirer en se ramenant tout de suite, pour la théorie des opérations $\Omega$, aux complexes de Kan, et aux chemins de type $\Delta_+$. Mais pour les composer \textit{fonctoriellement,} il me semble y avoir toujours le même problème - on trouvera des chemins de type $\Delta_+\Delta_+ = \boldsymbol{\Delta}_2^+$, et non pas $\Delta_+$, mais on peut alors composer les homotopies directes (revient à s'induire sur $\Delta_+\simeq\{0,2\}\subseteq\boldsymbol{\Delta}_2^+$). Or j'ai fait le formalisme des catégories de chemins sans me soucier de conditions du type Kan. Je présume que dans le cadre semi-simplicial, les $\uHom(\boldsymbol{\Delta}_1,X)$ ne donnent \textit{pas} le bon objet des chemins, pour $X$ pas de Kan. (Car sinon, il semblerait que pour $X=X_0$ à composantes \textit{finies,} le calcul des $\pi_i$ devrait être fini, ce qui n'est pourtant pas le cas, m'a-t-on toujours assuré \dots)

Je présume pourtant, quand on prend $\Sigma=\mathrm{Fl}(X)$, que les objets $\Ch_\infty(X;a,b)$ et $\Ch_\infty(\bar{X};a,b)$ sont isomorphes dans Hot, quand on immerge $X$ dans une enveloppe de Thomason-Kan $\bar{X}$.

\clearpage
\printedpage{81}
% head: A. Grothendieck : Dérivateurs, Ch. VII, Édition des Universités de Montpellier II et Paris VII
% foot: Version 22 décembre 2001

\msp{page 129}

\section*{11\quad Rapide retour sur les $(\mathcal{M},\Sigma)$, dans l'optique des dérivateurs}

Quand on s'intéresse au formalisme des $\Ch_\Sigma(X;a,b)$ en vue d'un \textit{dérivateur} $\mathbf{D}_{X,\Sigma}$ qu'on voudrait définir en termes de $X$, $\Sigma$, il semblerait que le prédérivateur qu'on obtient bel et bien ne mérite le nom de dérivateur que lorsqu'on suppose que pour toute flèche
\[ Y\longrightarrow Y' \]
dans Cat qui est une Hot-équivalence, les flèches correspondantes
\[ \mathrm{Hom}_{\mathrm{Ind\,Cat\,hot}}(Y',H(a,b)) \longrightarrow \mathrm{Hom}_{\mathrm{Ind\,Cat\,hot}}(Y,H(a,b)), \]
ou encore
\[ \varinjlim_n \mathrm{Hom}_{\mathrm{Cat\,hot}}(Y',\Ch_{n,\Sigma}(a,b)) \longrightarrow \varinjlim_n \mathrm{Hom}_{\mathrm{Cat\,hot}}(Y,\Ch_{n,\Sigma}(a,b)) \]
sont des bijections. J'ai regardé dans V ce que ça signifie pour l'objet $H(a,b) = (\Ch_{n,\Sigma}(a,b))_n$ de Ind\,Cat\,hot. C'est cette condition seulement qui assure, du point de vue des dérivateurs, une interprétation convaincante de $\Ch_{\infty,\Sigma}(a,b)$,
\[ \Ch_{\infty,\Sigma}(a,b) \;=\; \varinjlim_n \Ch_{n,\Sigma}(a,b) \]
en tant qu'objet de Hot, ou plus exactement de $\Hot^{\wedge}_{0}$, car on a alors pour $Y$ une catégorie \textit{finie}
\[ \underbrace{\mathrm{Hom}_{\mathrm{Cat\,hot}}(Y,H(a,b))}_{\simeq\,\mathrm{Hom}_{X(Y)\Sigma_Y^{-1}}(a_X,b_X)} \;\simeq\; \mathrm{Hom}_{\Hot}(Y,\Ch_{\infty,\Sigma}(a,b)), \]
\msp{page 130}
car je veux quand-même que le type dans Hot soit le bon, qu'il correspond bien à l'homotopie des espaces de chemins. Donc on voudrait qu'un morphisme
\[ X\xrightarrow{\;f\;}X' \]
dans (Cat), qui est une Hot-équivalence, définisse
\[ \Ch_\infty(X;a,b) \longrightarrow \Ch_\infty(X';f(a),f(b)), \]
qui soient eux aussi des Hot-équivalences. D'autre part, c'est évidemment compatible avec la brillante composition des chemins, que j'ai réussi à définir entre ces objets de Hot. Donc (sous réserve d'avoir vérifié ce que je viens de dire) la vérification de l'associativité se ramenerait au cas où $X$ est de Thomason-Kan. Mais dans ce cas, il y a des chances que, tout comme dans le cas semi-simplicial, les $\Ch_\infty(X;a,b)$ peuvent se définir (comme objets de (Hot)) comme étant simplement les chemins de type $\Delta_+$ de $a$ à $b$, qu'on compose alors de la façon triviale. Mais non, ces catégories sont discrètes! % [sic] ramenerait -> ramènerait, missing accent

Je renonce pour l'instant à poursuivre plus la question d'associativité. Demander à \textsc{Thomason} pour les propriétés du foncteur $\Hot\longrightarrow\Hot^{\wedge}_{0}$.

\clearpage
\printedpage{82}
% head: A. Grothendieck : Dérivateurs, Ch. VII, Édition des Universités de Montpellier II et Paris VII
% foot: Version 22 décembre 2001

\msp{page 131}
Cela signifie, à très peu de choses près, qu'on a une interprétation de $\Ch_{\infty,\Sigma}(a,b)$, en tant qu'objet de Hot, comme un `Hom externe' $\mathrm{Hom\,ext}(a,b)$ à valeurs dans Hot, défini sur la catégorie $\mathcal{M}=X$ à l'aide du seul ensemble admissible de flèches $\Sigma\subseteq\mathrm{Fl}(\mathcal{M})$, opération satisfaisant la condition
\[ \boxed{\mathrm{Hom}_{\Hot}(Y,\mathrm{Hom\,ext}(a,b)) \;\simeq\; \mathrm{Hom}_{\mathcal{M}(Y)\Sigma_Y^{-1}}(a_X,b_X)} \tag*{$(*)$} \]
(qui peut s'interpréter aussi, soit comme
\[ \mathrm{Hom}_{\mathcal{M}}(p_!(a_Y),b) \]
si $p_! : \underbrace{\mathcal{M}(Y)\Sigma_Y^{-1}}_{\mathbf{D}(Y)} \longrightarrow \mathcal{M}$ existe, soit comme
\[ \mathrm{Hom}_{\mathcal{M}}(a,p_*(b_Y)) \]
si $p_* : \underbrace{\mathcal{M}(Y)\Sigma_Y^{-1}}_{\mathbf{D}(Y)} \longrightarrow \underbrace{\mathcal{M}}_{=\,\mathbf{D}(*)}$ existe); avec le seul petit ennui que cette formule n'est valable que pour $Y$ \textit{finie.} Si on veut exprimer le second membre pour $Y$ variable quelconque dans Cat, il ne suffit pas d'un objet de Hot ($\mathrm{Hom\,ext}(a,b)$) dans le premier membre, mais il faut l'ind-objet $H(a,b) = (\Ch_n(a,b))_{n\in\mathbf{N}}$ dans Hot. Si donc on tient à travailler dans un `domaine' Diag du dérivateur $\mathbf{D}=\mathbf{D}_{\mathcal{M},\Sigma}$ aussi grand que possible, en se bornant à des catégories finies,
\msp{page 132}
et qu'on veut qu'une formule $(*)$ reste valable pour toute $Y$ dans le domaine, il faut `en toute rigueur' définir le $\mathrm{Hom\,ext}(a,b)$ non comme objet de Hot, mais de Ind\,Hot. Il se décrirait donc alors comme
\[ \mathrm{Hom\,ext}(a,b) = (\Ch_{n,\Sigma}(a,b))_n \qquad \text{(ind-objet dans Hot)}\ , \]
et l'objet $\Ch_{\infty,\Sigma}$ dans Hot en est une sorte de limite inductive.

Je me suis demandé si la condition préalable, savoir que les foncteurs $H(a,b)$ sur Cat transforment Hot-équivalences $Y'\longrightarrow Y$ en bijections, était bel et bien une (forte) restriction. Je me suis convaincu que c'est bien le cas, déjà dans le cas particulier où $\Sigma=\mathrm{Fl}(X)$. Il me semble que lorsque les composantes connexes de $\Omega(X,a)$ ne sont pas contractiles, ladite condition n'est pratiquement jamais satisfaite. J'ai pensé p.~ex. au cas où $X$ a le type d'homotopie (dans Hot) d'une 2-sphère, de sorte que son $\Omega(X,a)=\Ch_\infty(X,a,a)$ doit avoir un $\pi_1\simeq\mathbf{Z}$. Prenons $Y$ ayant le type dans Hot d'un cercle, disons
\[
\begin{tikzcd}
 & & a & & \\
Y\colon & 0 \arrow[ur] \arrow[dr] & & 0' \arrow[ul] \arrow[dl] & , \\
 & & b & &
\end{tikzcd}
\]

\clearpage
\printedpage{83}
% head: A. Grothendieck : Dérivateurs, Ch. VII, Édition des Universités de Montpellier II et Paris VII
% foot: Version 22 décembre 2001

\msp{page 133}
donc les foncteurs
\[ f : Y\longrightarrow \Ch_\infty(X;\bullet,\bullet) \overset{\text{déf}}{=} \Omega(X,\bullet) \]
correspondent à \textit{certains} `lacets' de $\Omega(X,\bullet)$ (mais lacets à origine non précisée), ils ont un type d'homotopie au sens topologique (`nombre de tours' autour de la sphère $X$) donné par un entier $n\in\mathbf{Z}$ (le `degré' du lacet $\lambda$). L'intuition évidente, c'est que l'ensemble des degrés $\delta$ qu'on peut obtenir ainsi, avec une catégorie \textit{finie} donnée $Y$ qui représente le cercle $S^1\in(\Hot)$, est fini - $|\delta|$ majoré sûrement par le nombre de ces flèches (donc au moins par 4 dans le cas envisagé, chaque flèche de $Y$ faisant faire \textit{au plus} un tour autour du pseudo-cercle $\Omega(S^2,\bullet)$) (\footnote[28]{probablement on ne trouve tout au plus que les degrés $\delta$ entre $-2$ et $+2$.}). D'autre part, quand on représente $S^1$ par des catégories $Y_n$ arbitrairement grandes (disons déduites des intervalles $\Phi_n=(\Delta_+\Delta_-)^n$ en identifiant la source et le but), on doit trouver des degrés arbitrairement élevés. Or toute $Y_n$ s'envoie dans $Y(=Y_2)$ par une flèche
\[ Y_n = Y' \longrightarrow Y = Y_2 \]
\msp{page 134}
qui est de degré 1, donc qui est une Hot-équivalence. Mais l'application correspondante
\[ \pi_0\uHom(Y,\Omega(X,\bullet)) \longrightarrow \pi_0\uHom(Y',\Omega(X,\bullet)) \]
n'est pas bijective. Ce genre d'argument me semble devoir marcher chaque fois que l'ensemble des groupes d'homotopie $\pi_j(X,x)$, pour $j\geq 2$ (donc l'ensemble des groupes d'homotopie $\pi_i$ de $\Omega(X,x)$, pour $i\geq 1$) est \textit{infini,} en prenant pour $Y$, $Y'$ différents représentants $Y_n$ de la sphère $S^i$, disons les $Y_n$ déduites de $\Phi_n$ $(n\geq 2)$ en prenant $\Phi_n^i = \Phi_n\times\Phi_n\times\cdots\times\Phi_n$ (qui est une sorte de $i$-cube combinatoire), et en identifiant son `bord' à un seul point. Je présume que le nombre d'élements dans $\pi_i(\Omega(X,x))$ qu'on peut obtenir à partir de flèches $Y_n\longrightarrow\Omega(X,x)$ est limité par la taille de $Y_n$, mais on peut sûrement les obtenir chacun, en prenant $n$ assez grand. Si donc $\xi\in\pi_i(\Omega)$ échappe à $Y_{n_0}$ ($n_0=0$, disons), on prend un $n'$ assez grand pour que $\xi$ n'échappe pas à $Y_{n'}$, et on choisit % [sic] d'élements -> d'éléments, missing accent
\[ Y_{n'} \longrightarrow Y_n\ , \]
\texttt{[}$n=n_0$\texttt{?]}

\msp{page 135}
qui soit une Hot-équivalence (p.~ex. à l'aide d'une rétraction de $\Phi_{n'}$ sur $\Phi_n$, compatible avec sources et buts).

Ceci attire fortement l'attention sur le fait que si $Y$, $C$ sont deux catégories, la catégorie $\uHom(Y,C)$ \textit{n'est pas} un représentant de $\uHom_{\Hot}(\mathrm{hot}(Y),\mathrm{hot}(C))$, on a seulement une flèche canonique
\[ \mathrm{hot}(\uHom(Y,C)) \longrightarrow \uHom_{\Hot}(\mathrm{hot}(Y),\mathrm{hot}(C)), \tag*{$(*)$} \]

% ---- from Der-ChVII-456_p086-090.tex
\printedpage{84}
% head: A. Grothendieck : Dérivateurs, Ch. VII, Édition des Universités de Montpellier II et Paris VII
% foot: Version 22 décembre 2001

i.e.
\[
\underbrace{\hot(\uHom(Y,C))\times\hot(Y)}_{\hot(\uHom(Y,C)\times Y)}\longrightarrow\hot(C),
\]
laquelle flèche est simplement déduite de
\[
\uHom(Y,C)\times Y\longrightarrow C.
\]
Mais la flèche $(*)$ n'a aucune envie d'être \texttt{[un]} isomorphisme. C'est parce qu'elle ne l'est déjà pas dans le cas particulier $Y=\mathbf{\Delta}^1$, essentiellement, qui nous oblige, pour faire une théorie des chemins dans les catégories qui soit topologiquement raisonnable, de ne pas nous borner à $\mathbf{\Delta}^1$ comme type de chemins, mais à prendre des intervalles composés $\tau\in\mathrm{Ty}$.

\msp{page 136}

Par contre, $(*)$ est \texttt{[un]} isomorphisme quand $C$ est une catégorie de Thomason-Kan. Il suffit de voir que pour tout $Z$ dans (Cat), le Hom de $\hot(Z)$ dans les deux membres est isomorphe par la flèche déduite de $(*)$. Or si $C$ est de Thomason-Kan, de même $\uHom(Y,C)$ (immédiat sur la définition de Thomason-Kan en termes de prolongements de
\begin{tikzcd}[baseline=(\tikzcdmatrixname-1-1.base), row sep=small]
A_0 \arrow[r, hook, "i"] \arrow[dr] & A \arrow[d, dotted] \\
& C
\end{tikzcd}
$\ldots$(\footnote[29]{un peu vif! Il faut savoir que les applications cofibrantes $i$ de \textsc{Thomason} le restent après multiplication cartésienne par $Y$ \ldots})), donc
\[
\mathrm{Hom}_{\mathrm{Hot}}(\hot(Z),\hot(\uHom(Y,C)))\xleftarrow{\ \sim\ }\pi_0(\underbrace{\uHom_{\mathrm{Cat}}(Z,\uHom(Y,C))}_{\uHom_{\mathrm{Cat}}(Z\times Y,C)})\simeq\pi_0\uHom(Z\times Y,C), % Reading resolved against the supplied PDF; see corrections.csv.
\]
et
\[
\begin{aligned}
&\mathrm{Hom}_{\mathrm{Hot}}(\hot(Z),\uHom_{\mathrm{Hot}}(\hot(Y),\hot(C)))\simeq\mathrm{Hom}_{\mathrm{Hot}}(\underbrace{\hot(Z)\times\hot(Y)}_{\hot(Z\times Y)},\hot(C))\\
&\qquad\underbrace{\xleftarrow{\ \sim\ }}_{\text{car $C$ est de Thomason-Kan}}\qquad\quad\pi_0(\uHom_{\mathrm{Cat}}(Z\times Y,C)),
\end{aligned}
\]
O.K.

On voit donc que les catégories $\uCh_{\infty,\Sigma}(X;a,b)$ n'ont aucune tendance à être des catégories de Thomason-Kan. Pourtant, son type dans Hot doit être `le bon', sans avoir d'abord à remplacer $X$ par une enveloppe de Thomason-Kan $\bar{X}$.

\printedpage{85}
% head: A. Grothendieck : Dérivateurs, Ch. VII, Édition des Universités de Montpellier II et Paris VII
% foot: Version 22 décembre 2001

\msp{page 137}

\section*{12\quad Retour sur les `catégories de chemins' dans \textmd{Cat}, comme catégorie de modèles pour \textmd{Hot}}

Je vois exactement deux applications intéressantes du formalisme des chemins dans les catégories. L'une, qui m'a surtout forcé la main, est pour pouvoir construire des Hom\,ext, Hom \textit{externes} à valeurs dans Hot, pour tout prédérivateur $\mathbf{D}$ défini par un couple $(\mathcal{M},\Sigma)$ (et je ne connais pas de prédérivateur qui ne soit défini ainsi).

L'autre application est pour avoir dans (Cat) une \textit{construction fonctorielle,} en termes d'un objet $X$ quelconque ou d'une flèche $X\longrightarrow Y$, de l'analogue du formalisme de \textsc{Cartan-Serre} sur les espaces de chemins. Dans ce cadre, il est entendu qu'on prend toujours $\Sigma_X=\mathrm{Fl}(X)$, ensemble de \textit{toutes} les flèches de $X$. Pour construire fonctoriellement en $X$ dans (Cat)
\[
\begin{tikzcd}[row sep=huge]
\uCh_\infty(X) \arrow[d] \\
X\times X,
\end{tikzcd}
\]
il faut avoir fixé un système projectif
\[
\underline{t}=(t_n)_{n\in\mathbf{N}}
\]
\msp{page 138}
dans \underline{Ty} qui soit cofinal dans \underline{Ty}, ce qui signifie simplement que parmi les $t_n$ il y en a de longueur utile arbitrairement grande (et comme
\[
n\longmapsto\mathrm{long.ut.}(t_n)
\]
est croissante, cela implique que long.ut.$(t_n)$ tend vers $+\infty$ pour $n\longrightarrow+\infty$). Les faits principaux à vérifier ensuite (si faire se peut) sont les propriétés essentielles des flèches dans le diagramme
\printedpage{86}
% head: A. Grothendieck : Dérivateurs, Ch. VII, Édition des Universités de Montpellier II et Paris VII
% foot: Version 22 décembre 2001
\[
\begin{tikzcd}[column sep=large, row sep=large]
& X \arrow[d, "\tau"] & \\
& \uCh_\infty(X) \arrow[d, "p"] \arrow[ddl, "\sigma"'] \arrow[ddr, "\beta"] & \\
& X\times X \arrow[dl, "\mathrm{pr}_1"] \arrow[dr, "\mathrm{pr}_2"'] & \\
X & & X\ ,
\end{tikzcd}
\qquad\qquad
\left|\begin{array}{l}
\sigma\tau=\beta\tau=\mathrm{id}_X\\
p\tau=\mathrm{diag}_X
\end{array}\right.
\]
et plus généralement, pour une flèche
\[
X'\xrightarrow{\ f\ }X
\]
dans Cat, des flèches dans le diagramme correspondant
\[
\begin{tikzcd}
& \overline{X'}=C(f) \arrow[dl] \arrow[dr, "g"] \arrow[ddrr, bend left=30, "\bar{f}"] & & \\
X' \arrow[ur, dotted, bend left=40, "\alpha"] \arrow[dr, "f"'] & |[draw, rounded corners]| \text{cart.} & \uCh_\infty(X) \arrow[dl, "\sigma"] \arrow[dr, "\beta"] & \\
& X & & X,
\end{tikzcd}
\qquad
\bar{f}\alpha=(\beta g)\alpha=f,
\]
et plus particulièrement de $\beta g$. La flèche
\msp{page 139}
en pointillés $\alpha$ est la section de $\overline{X'}$ sur $X'$ déduite de la section $\tau$ de $\uCh_\infty(X)$ sur $X$ qui réalise la décomposition de $f$ en
\[
X'\xrightarrow{\ \alpha\ }\overline{X'}\xrightarrow{\ \bar{f}\ }X,
\]
où $\bar{f}$ devrait être une flèche \textit{`fibrante'} et une Hot-équivalence, $\alpha$ une flèche \textit{`cofibrante'} au sens de \textsc{Thomason}. (Mais s'il en est ainsi, est à voir!) Voici pour le moment les propriétés essentielles que je voudrais vérifier.

\begin{enumerate}[label={}, leftmargin=*]
\item[1$^o$)] $p$ est lisse et propre, et est une HOT-fibration.

\item[2$^o$)] Donc il en est de même de $\sigma=\mathrm{pr}_1p$ et $\beta=\mathrm{pr}_2p$. De plus, ce sont des HOT-fibrations \textit{triviales,} i.e. à fibres asphériques.

\printedpage{87}
% head: A. Grothendieck : Dérivateurs, Ch. VII, Édition des Universités de Montpellier II et Paris VII
% foot: Version 22 décembre 2001

\item[3$^o$)] Plus généralement (pour ce qui concerne $\beta$), $\overline{X'}\xrightarrow{\ \bar{f}\ }X$ devrait encore être propre et lisse et une HOT-fibration, mais bien sûr en général pas triviale. (Elle devrait être triviale si et seulement si $f$ est une Hot-équivalence.) (\footnote[30]{résulte formellement de 1$^o$ et 2$^o$, cf. page 170.}).
\end{enumerate}

Un mot sur la `dualité'. Les énoncés sur $\sigma$ (resp. $\beta$) se traduisent en énoncés sur $\beta$ (resp. $\sigma$) par passage de $\underline{t}$ à $\underline{t}^o=(t^o_n)$, où $\tau\longmapsto\tau^o$ est le foncteur dans \underline{Ty} `passage à l'intervalle opposé'. On a en effet % Reading resolved against the supplied PDF; see corrections.csv.

\msp{page 140}

un isomorphisme fonctoriel en $X$ (passage au chemin opposé)
\[
\boxed{\begin{array}{ccc}
c & \longmapsto & c^o\\[4pt]
\uCh^{\underline{t}}_\infty(X) & \longrightarrow & \uCh^{\underline{t}^o}_\infty(X)
\end{array}}
\qquad\qquad
\text{(\footnotemark[31])}
\]
\footnotetext[31]{cf. page 143 pour le passage de $\underline{t}$ à $\underline{t}^-$.}
qui `échange source et but'. Bien sûr, on a intérêt à prendre $\underline{t}$ tel qu'on ait $\underline{t}=\underline{t}^o$ (c'est ce que j'avais fait, par un sain instinct, dès le début dans V), de sorte qu'on a un foncteur involutif canonique
\[
\begin{array}{ccc}
c & \longmapsto & c^o\\[4pt]
\uCh_\infty(X) & \longrightarrow & \uCh_\infty(X)
\end{array}
\]
(je laisse tomber à nouveau l'exposant $\underline{t}$), qui `échange source et but'. Donc on trouve, à isomorphisme canonique près, la `même' factorisation de $f:X'\longrightarrow X$ en $\bar{f}\circ\alpha$, si on fait comme dans page 38 pour définir $\overline{X'}$ et $\alpha$, ou si dans ce diagramme on intervertit les rôles de $\sigma$ et de $\beta$. % [sic] page 38 -> page 138, the diagram defining \overline{X'} and \alpha stands on manuscript page 138 (apparently)

Enfin :

\begin{enumerate}[label={}, leftmargin=*]
\item[4$^o$)] Si
\[
c:a\longrightarrow b
\]
est un chemin quelconque dans $X$, on voudrait en conclure un isommorphisme dans Hot (de `transport par $c$') % [sic] isommorphisme -> isomorphisme, doubled m
\[
\overline{X'}_a\xrightarrow[\sim]{\bar{c}}\overline{X'}_b,
\]
avec bien sûr transitivité pour la composition,

\msp{page 141}

ce qui donnera aussi
\[
\overline{c^{\mathrm{opp}}}=\overline{c}^{-1}
\]
(en admettant que les chemins triviaux $a\xrightarrow{\mathrm{id}}a$ donnent l'identité $\overline{X'}_a\xrightarrow[\sim]{\mathrm{id}}\overline{X'}_a$ !) (\footnotemark[32]\textsuperscript{,}\footnotemark[33]). Bien sûr, $\overline{c}$ ne doit dépendre que de la \textit{classe d'homotopie} de $\overline{c}$ (dans % [sic] de $\overline{c}$ (dans -> de $c$ (dans, the homotopy class is apparently that of the path c
\footnotetext[32]{Surtout pas se borner au cas où $c$ serait un des $t_n$!}
\footnotetext[33]{décrire quand-même l'effet sur les $\pi_0$!}
\printedpage{88}
% head: A. Grothendieck : Dérivateurs, Ch. VII, Édition des Universités de Montpellier II et Paris VII
% foot: Version 22 décembre 2001
$\uCh_\tau(X;a,b)$, ou même dans $C=\uCh(X;a,b)$, catégorie de tous les chemins dans $X$, de $a$ vers $b$, de type non précisé). En d'autres termes, on doit avoir un foncteur
\[
\begin{array}{rcl}
\Pi(X) & \longrightarrow & \mathrm{Hot}\\[4pt]
a & \longmapsto & \hot(\overline{X'}_a)
\end{array}
\]
du groupoïde fondamental de $X$ dans Hot. À vrai dire, ceci n'a rien à voir avec la situation particulière étudiée ici - ça doit marcher chaque fois qu'on a une HOT-fibration, ici $\overline{X'}\longrightarrow X$, mais aussi $\uCh_\infty(X)\longrightarrow X\times X$ (\footnote[34]{Ici aussi, il faudrait décrire l'effet du transport $\overline{c}$ sur les $\pi_0$ des fibres.}). Et pour une $\mathbf{D}$-fibration $X'\longrightarrow X$ ($\mathbf{D}$ un dérivateur quelconque), et $a\in\mathbf{D}(e)$, on doit avoir de même deux foncteurs
\[
\Pi(X)\rightrightarrows\mathbf{D}(e)
\]
correspondant à
\[
\left\{\begin{array}{lcl}
x & \longmapsto & H^\bullet_{\mathbf{D}}(X'_x,\bullet)\\[4pt]
x & \longmapsto & H_\bullet^{\mathbf{D}}(X'_x,\bullet).
\end{array}\right.
\]

\msp{page 142}

Le cas qui nous occupe ici serait celui où $\mathbf{D}=\mathrm{HOT}$, où on prend $H^{\mathbf{D}}_\bullet$, et $a=$ objet final de $\mathbf{D}(e)=\mathrm{Hot}$.)

\item[5$^o$)] Si 1$^o$, 2$^o$, 3$^o$ sont vérifiés, on est moralement sûr que la construction des $\uCh_\infty(X)$ est bel et bien l'équivalent exact du formalisme de \textsc{Cartan-Serre}. Pour en avoir le c\oe ur totalement net, i.e. raccrocher explicitement cette théorie aux constructions classiques, il faudrait prouver un \textit{énoncé de comparaison} avec la théorie semi-simpliciale, qui dirait essentiellement ceci : Si on désigne, pour $X$ dans Cat, par $\tilde{X}$ l'objet associé dans $\mathbf{\Delta}^\wedge$ (le `nerf' de $X$), alors la flèche
\[
\uCh_\infty(X)\tilde{}\longrightarrow\underbrace{(X\times X)\tilde{}}_{\simeq\,\tilde{X}\times\tilde{X}}
\]
dans $\mathbf{\Delta}^\wedge/(X\times X)\tilde{}$, est $W$-équivalente sur $(X\times X)\tilde{}$ à la flèche
\[
\uCh(\tilde{X})\longrightarrow\tilde{X}\times\tilde{X}
\]
de la théorie semi-simpliciale des chemins (à supposer qu'il y ait une telle flèche classique définie pour \textit{tout} objet $X$ de $\mathbf{\Delta}^\wedge$). Ou qu'on ait tout au moins des isomorphismes canoniques dans Hot
\[
\hot_{\mathrm{Cat}}(\uCh_\infty(X;a,b))\ \simeq\ \hot_{\mathbf{\Delta}^\wedge}(\uCh(\tilde{X};\tilde{a},\tilde{b})).
\]
\end{enumerate}

% ---- from Der-ChVII-456_p091-095.tex
\printedpage{89}
% head: A. Grothendieck : Dérivateurs, Ch. VII, Édition des Universités de Montpellier II et Paris VII
% foot: Version 22 décembre 2001

\msp{page 143}
Dans $1^o$ et $3^o$, on veut prouver que certaines flèches ($p$ et $\bar{f}$) sont à la fois \textit{propre} et \textit{lisse} - ce qui est vraiment idéal ! Je dis qu'il suffit à chaque fois de prouver l'un des deux, que l'autre s'en déduit (comme `cas particulier') par dualité. Pour ceci, notons
\[
\tau \longmapsto \tau^{-}
\]
l'automorphisme (\textit{pas anti}-automorphisme) de $\Ty$ qui consiste à changer de signe toutes les flèches de transition - c'est donc, au niveau des objets, l'unique automorphisme du monoïde libre Ty, qui échange $\Delta_+$ et $\Delta_-$. Au système projectif $\underline{t}$ dans $\Ty$ correspond le système projectif $\underline{t}^-$. Ceci dit, on a un isomorphisme de catégories
\[
\boxed{\Ch^{\underline{t}}_{\infty}(X^o) \ \simeq\ (\Ch^{\underline{t}^-}_{\infty}(X))^o} \tag*{$(^{35})$} % Reading resolved against the supplied PDF; see corrections.csv.
\]
\footnotetext[35]{cf. page 140 pour le passage de $\underline{t}$ à $\underline{t}^o$.}
et le diagramme commutatif
\[
\begin{tikzcd}
\Ch^{\underline{t}}_{\infty}(X^o) \arrow[r, "\sim"] \arrow[d, "p^{\underline{t}}_{X^o}"] & \Ch^{\underline{t}^-}_{\infty}(X)^o \arrow[d, "(p^{\underline{t}^-}_{X})^o"] \\
X^o \times X^o \arrow[r, equal] & (X \times X)^o.
\end{tikzcd}
\]
Ainsi, dire que $p^{\underline{t}}_{X}$ \texttt{[plutôt} $p^{\underline{t}^-}_{X}$\texttt{]} est lisse, i.e. que $(p^{\underline{t}}_{X})^o$ \texttt{[plutôt} $(p^{\underline{t}^-}_{X})^o$\texttt{]} est propre, équivaut à dire que $p^{\underline{t}}_{X^o}$ est propre.

\msp{page 144}
Il nous reste donc, à ce sujet, que deux faits à vérifier : la lissité de $p : \Ch_{\infty}(X) \longrightarrow X \times X$, et celle de $\bar{f} : \overline{X'} \longrightarrow X$.

a) Lissité de $p : \Ch_{\infty}(X) \longrightarrow X \times X$. On se donne une flèche
\[
w : (a,b) \longrightarrow (a',b'), \qquad \text{i.e.} \quad \left\{ \begin{array}{lccc} u : & a & \longrightarrow & a' \\[1ex] v : & b & \longrightarrow & b', \end{array} \right.
\]
sur la base, un objet $c'$ au dessus du but, i.e.
\[
c' \in \Ch_{\infty}(X; a', b'),
\]
et on regarde la catégorie des flèches dans $\Ch_{\infty}(X)$, de but $c'$, qui relèvent $w$,
\[
\begin{tikzcd}
\underset{\text{\scriptsize variable}}{c} \arrow[r, "\varphi"] \arrow[d, dash] & \underset{\text{\scriptsize fixe}}{c'} \arrow[d, dash] \\
(a,b) \arrow[r, "w"] & (a',b'),
\end{tikzcd}
\]
\printedpage{90}
% head: A. Grothendieck : Dérivateurs, Ch. VII, Édition des Universités de Montpellier II et Paris VII
% foot: Version 22 décembre 2001
i.e. la catégorie des couples $(c,\varphi)$, où
\[
\left\{ \begin{array}{l}
c \in \Ch_{\infty}(X; a, b) \\[1ex]
\varphi : c \longrightarrow c' \ \ (\in \Ch_{\infty}(X; a', b')) \\[1ex]
\text{une flèche induisant } \begin{array}[t]{l} u \ \text{ sur les sources} \\ v \ \text{ sur les buts,} \end{array}
\end{array} \right.
\]
\[
\begin{tikzcd}
a \arrow[rrrr, dotted, dash, "c"] \arrow[d] & {} \arrow[d, dotted] & {} \arrow[d, dotted] & {} \arrow[d, dotted] & b \arrow[d] \\
a' \arrow[rrrr, dash, "c'"] & {} & {} & {} & b'.
\end{tikzcd}
\]

\msp{page 145}
\begin{flushright}
{\footnotesize 9.11.90}
\end{flushright}

Pour la question précédente, relative aux flèches dans le diagramme
\[
\begin{tikzcd}
& X \arrow[d, "i"] & \\
& \Ch_{\infty}(X) \arrow[d, "p"] \arrow[ddl, "\sigma"'] \arrow[ddr, "\beta"] & \\
& X \times X \arrow[dl, "\mathrm{pr}_1"] \arrow[dr, "\mathrm{pr}_2"'] & \\
X & & X,
\end{tikzcd}
\]
il faut préciser absolument la déscription de $\Ch_{\infty}(X)$ \textit{dans} Cat, donc il faut préciser le système projectif $\underline{t} = (t_n)_{n \in \mathbf{N}}$. Nous allons prendre en effet % [sic] déscription -> description, misspelling
\[
t_n = (\underbrace{\Delta_+ \Delta_-}_{\tau_0})^{2n},
\]
avec les morphismes de transition décrits par ailleurs
\[
\underbrace{t_{n+p}}_{=\tau_0^{2n+2p}} \ \longrightarrow \ \underbrace{t_n}_{\tau_0^{n}}, % [sic] \tau_0^{n} -> \tau_0^{2n}, since t_n = (\Delta_+\Delta_-)^{2n} = \tau_0^{2n}
\]
\printedpage{91}
% head: A. Grothendieck : Dérivateurs, Ch. VII, Édition des Universités de Montpellier II et Paris VII
% foot: Version 22 décembre 2001
étant donnés par $\delta_p * \mathrm{id}_{\tau_o^{2n}} * \delta_p$, où $\delta_p : \tau_0^p \longrightarrow 1$ est la dégénérescence de $\tau_0^p$. Ainsi on peut considérer $\Ch_{\infty}(X)$ comme étant la catégorie des chemins biïnfinis dans $X$ % [sic] \tau_o^{2n} -> \tau_0^{2n}, letter o printed in the index instead of 0
\[
\dcfit{\cdots \longrightarrow a_{-5} \longleftarrow \overbrace{a_{-4} \longrightarrow a_{-3} \xleftarrow{u_{-2}} \overbrace{a_{-2} \xrightarrow{u_{-1}} a_{-1} \xleftarrow{u_0} \overbrace{a_0}^{\Ch_0} \xrightarrow{u_1} a_1 \xleftarrow{u_2} a_2}^{\Ch_1} \xrightarrow{u_3} a_3 \longleftarrow a_4}^{\Ch_2} \longrightarrow a_5 \longleftarrow \cdots,}
\]
avec les $u_i$ des identités pour $|i|$ assez grand.

La sous-catégorie (pleine)
\[
\Ch^{W}_{t_n}(X) \ = \ \Ch^{W}_{(\Delta_+\Delta_-)^{2n}}
\]
s'identifie à celle pour laquelle on a
\begin{quote}
$u_i$ une identité pour $i \notin [-2i+1, 2i]$, i.e. les seuls $u_i$ `significatifs' sont les $4i$ flèches $u_{-2i+1}$, $u_{-2i+2}$, $\ldots$, $u_{2i}$. % [sic] i (in [-2i+1,2i], 4i, u_{-2i+1}, u_{-2i+2}, u_{2i}) -> n, the index of t_n is n
\end{quote}

Mais il sera plus commode de regarder $\Ch_{\infty}(X)$ comme la limite inductive d'une
\msp{page 146}
famille à deux indices de sous-catégories pleines
\[
\Ch_{n,m}(X), \qquad n, m \in \mathbf{Z},
\]
en définissant cette catégorie comme \textit{vide} si $n > m$, et pour $n \leq m$ comme formée des objets
\[
\left\{ \begin{array}{l}
(a_i, u_i)_{i \in \mathbf{Z}} \in \Ch^{W}_{n,m}(X) \iff u_i \text{ une identité pour } i \notin [n, m-1] \\[1ex]
\overset{\mathrm{id}}{\cdots}\, a_{n-1} \dashto{\mathrm{id}} \underbrace{a_n \dashto{u_n} a_{n+1} \cdots \dashto{} a_{m-1} \dashto{u_{m-1}} a_m}_{\text{\scriptsize `partie utile' du chemin}} \dashto{\mathrm{id}} a_{m+1} \overset{\mathrm{id}}{\cdots}\,.
\end{array} \right.
\]
Donc pour $n = m$, on trouve la sous-catégorie pleine isomorphe à $X$, formée des chemins infinis `totalement dégénérés', i.e. tels que \textit{toutes} les $u_i$ soient des identités. (Cette sous-catégorie de $\Ch_{\infty}$ ne dépend donc pas du choix de $n = m$.)
\[
\xrightarrow{\mathrm{id}} \underbrace{a_{-1}}_{=\, a} \xleftarrow{\mathrm{id}} \underbrace{a_0}_{=\, a} \xrightarrow{\mathrm{id}} \underbrace{a_1}_{=\, a} \xleftarrow{\mathrm{id}}
\]
Tous les sommets d'un tel chemin sont égaux à un même objet de $X$, et les flèches de transition entre eux sont identiques.

La structure de $\Ch^{W}_{n,m}(X) \ \simeq \ \Ch^{W}_{\tau_{n,m}}(X)$ ne dépend que de la longueur $m-n$ de $\tau_{n,m}$ et de la \textit{parité} de $n$. Si $n$ est pair, la direction initiale de $\tau_{n,m}$ est positive, et
\[
\tau_{n,m} = \underbrace{\Delta_+ \Delta_- \cdots}_{m-n \text{ facteurs en alternance}} = \left\{ \begin{array}{ll}
(\Delta_+ \Delta_-)^{\frac{m-n}{2}} & \text{si } m \text{ pair} \\[1ex]
(\Delta_+ \Delta_-)^{\frac{m-n}{2}-1} \Delta_+ & \text{si } m \text{ impair}
\end{array} \right.
\]
\texttt{[plutôt} $\frac{m-n-1}{2}$ \texttt{au lieu de} $\frac{m-n}{2} - 1$\texttt{]}, tandis que pour $n$ impair c'est juste l'inverse
\printedpage{92}
% head: A. Grothendieck : Dérivateurs, Ch. VII, Édition des Universités de Montpellier II et Paris VII
% foot: Version 22 décembre 2001
\msp{page 147}
\[
\tau_{n,m} = \underbrace{\Delta_- \Delta_+ \cdots}_{m-n \text{ facteurs alternés}} = \left\{ \begin{array}{ll}
(\Delta_- \Delta_+)^{\frac{m-n}{2}} & \text{si } m \text{ impair} \\[1ex]
(\Delta_- \Delta_+)^{\frac{m-n}{2}-1} \Delta_- & \text{si } m \text{ pair}
\end{array} \right.
\]
\texttt{[plutôt} $\frac{m-n-1}{2}$ \texttt{au lieu de} $\frac{m-n}{2} - 1$\texttt{]}. La direction initiale est conforme à la parité de $n$, i.e. égale à $(-1)^n$, tandis que la direction terminale est opposée à la parité de $m$, i.e. égale à $(-1)^m$ \texttt{[plutôt} $(-1)^{m+1}$\texttt{]}.

Je note que l'on a $\tau_0 = (\tau_0)^o$, d'où
\[
\underline{t} = \underline{t}^o, \qquad \text{mais } \underline{t} \neq \underline{t}^-,
\]
et pourtant on a des isomorphismes fonctoriels en $X$
\[
\Ch_{\underline{t}}(X) \ \simeq \ \Ch_{\underline{t}^-}(X)
\]
respectant source et but, à cause d'un isomorphisme
\[
\underline{t} \ \simeq \ \underline{t}^- \qquad \text{dans Pro}(\Ty).
\]
On peut dire p. ex. que le \textit{foncteur translation} $T^p$ d'amplitude $p$ (pour $p \in \mathbf{Z}$) est défini sur $\Ch^{W}_{\underline{t}}(X)$, par
\[
\left\{ \begin{array}{l}
T^p(((a_i),(u_i))_{i \in \mathbf{Z}}) = ((b_j),(v_j))_{j \in \mathbf{Z}} \\[1ex]
\text{avec} \left\{ \begin{array}{lcl} b_j & = & a_{j-p} \\[1ex] v_j & = & u_{j-p} \end{array} \right.
\end{array} \right.
\]
et que
\[
\begin{aligned}
T^p &: \left\{ \begin{array}{lcl} \Ch_{t}(X) & \xrightarrow{\sim} & \Ch_{-t}(X) \\[1ex] \Ch_{-t}(X) & \xrightarrow{\sim} & \Ch_{t}(X) \end{array} \right\} \ \text{si } p \text{ impair} \\
T^p &: \left\{ \begin{array}{lcl} \Ch_{t}(X) & \xrightarrow{\sim} & \Ch_{t}(X) \\[1ex] \Ch_{-t}(X) & \xrightarrow{\sim} & \Ch_{-t}(X) \end{array} \right\} \ \text{si } p \text{ pair.}
\end{aligned} % Reading resolved against the supplied PDF; see corrections.csv.
\]

\msp{page 148}
En d'autres termes, on peut formellement regarder la catégorie somme
\[
\Ch_{\pm}(X) = \Ch_{+}(X) \amalg \Ch_{-}(X)
\]
et l'opération translation $T$ comme un automorphisme qui échange les deux morceaux, et $T^p$ n'est autre que son itéré $p$-ème, qui échange les morceaux $+$ et $-$ si $p$ impair, qui les stabilise si $p$ pair.

\printedpage{93}
% head: A. Grothendieck : Dérivateurs, Ch. VII, Édition des Universités de Montpellier II et Paris VII
% foot: Version 22 décembre 2001

Il y a aussi l'automorphisme involutif de passage au chemin opposé, qui échange source et but, et que je définis (conformément à l'isomorphisme canonique identique
\[
\underline{t} \xrightarrow{\sim} \underline{t}^o \quad )
\]
comme la symétrie $S$
\[
\left\{ \begin{array}{l}
S((a_i),(u_i)) = ((b_j),(u_j)) \ \ \text{avec} \\[1ex]
\begin{array}{lcl} b_j & = & a_{-j} \\[1ex] u_j & = & u_{-j+1}. \end{array}
\end{array} \right.
\]
Si je le regarde comme opérant sur $\Ch^{W}_{\pm}(X)$, il est lié au foncteur $T$ par
\[
\left\{ \begin{array}{rcl}
S^2 & = & \mathrm{id} \\[1ex]
STS^{-1} & = & T^{-1},
\end{array} \right.
\]
de sorte que c'est le groupe semidirect de $\mathbf{Z}/2\mathbf{Z}$ par $\mathbf{Z}$ (groupe affine de dimension 1 de l'anneau $\mathbf{Z}$) qui opère sur $\Ch_{\pm}$.

\msp{page 149}
On est embarassé de trouver un automorphisme $T^2$ de la situation $\Ch_{\infty}(X) \longrightarrow X \times X$, sans qu'il soit évident que dans Hot, ce foncteur induise l'identité! Je ne vois aucun moyen de relier un chemin $c$ et $T^2 c$ par une flèche dans $\Ch_{\infty}$, ni même pas un chemin `canonique'. Mais je me dis que cet automorphisme est décrit par un automorphisme, non pas du système projectif % [sic] embarassé -> embarrassé, misspelling
\[
(\tau_{n,m})_{n,m \in \mathbf{Z}^o \times \mathbf{Z},\ n \leq m}
\]
(où $\mathbf{Z}^o$ désigne l'ensemble ordonné \textit{opposé} de $\mathbf{Z}$), mais du \textit{pro-objet} de $\Ty$ qu'il définit (qui est bien sûr `le même' pour tous les systèmes projectifs cofinaux à gauche dans $\Ty$). Il en résulte déjà que $T^2$ induit l'identité sur l'ind-objet de Cat\,hot, $\underline{\underline{\mathrm{Ch}}}_{\infty}(X)$, et a fortiori sur l'ind-objet correspondant de Hot, ou de $\mathrm{Hot}^{\wedge}_{0}$, ou enfin pour l'objet de $\mathrm{Hot}^{\wedge}_{0}$ lui-même. De même pour les $\Ch_{\infty}(X; a, b)$. Il s'ensuit, revenant à $\Ch_{\infty}(X)$ lui-même et aux $\Ch_{\infty}(X; a, b)$, que l'automorphisme induit par $T^2$ est l'identité dans $\mathrm{Hot}^{\wedge}_{0}$. J'espère que $\mathrm{Hot} \longrightarrow \mathrm{Hot}^{\wedge}_{0}$ est
\msp{page 150}
tout au moins fidèle, de sorte que ça prouvera que $T^2$ induit l'identité dans Hot également.

Je reviens aux questions $1^o$ à $4^o$ d'il y a quelques jours.

Pour $1^o$), il faut prouver que
\[
p : \Ch_{\infty}(X) \longrightarrow X \times X
\]
est lisse et propre et une Hot-fibration. Par dualité, $p$ propre résultera de $p$ lisse. Mais $\Ch_{\infty}(X)$ est $\varinjlim$ des $\Ch_{n,m}(X)$, avec $n \longrightarrow -\infty$, $m \longrightarrow +\infty$, et où on peut se borner à $n$ et $m$ \textit{pairs.} Or on voit que pour $n$, $m$ pairs, i.e. $\tau_{n,m}$ à flèche initiale directe, à flèche finale

% ---- from Der-ChVII-456_p096-100.tex
\printedpage{94}
% head: A. Grothendieck : Dérivateurs, Ch. VII, Édition des Universités de Montpellier II et Paris VII
rétrograde, $\Ch_{n,m}(X) \simeq \Ch_{\tau_{n,m}}(X)$ est Cat-\textit{fibrée} sur $X \times X$, donc lisse (i.e. HOT-lisse), et cela passe aux $\varinjlim$ (\footnote[36]{\textbf{NB} $\Ch_\infty(X)$ n'est ni fibrée ni cofibrée sur $X \times X$ - car les foncteurs d'inclusion entre les $\Ch_{n,m}$ ne sont pas cartésiens resp. cocartésiens.}).

Ainsi $\Ch_\infty(X) \xrightarrow{p} X \times X$ est à la fois propre et lisse (\footnote[37]{12.11. Mais ça implique déjà que c'est une fibration !}). Dire que c'est une HOT-fibration, revient p. ex. à dire que pour
\[
(a,b) \xrightarrow{(u,v)} (a',b')
\]
dans $X \times X$, la flèche correspondante
\[
\Ch_\infty(X)/(a,b) \longrightarrow \Ch_\infty(X)/(a',b')
\]
est une Hot-équivalence. Or, comme $p$ est propre, les inclusions des fibres en $(a,b)$ et en $(a',b')$ dans l'un et l'autre membre sont des Hot-équivalences :
\msp{page 151}
\[
\begin{tikzcd}[row sep=large]
\Ch_\infty(X)/(a,b) \arrow[r, "\rho"] & \Ch_\infty(X)/(a',b') \\
\Ch_\infty(X;a,b) \arrow[u, hook, "\wr"] \arrow[r, dotted, "\varphi"] & \Ch_\infty(X;a',b') \arrow[u, hook, "\wr"]
\end{tikzcd}
\tag*{$(*)$}
\]
Ainsi on trouve une flèche en pointillés dans Hot, dont il faut prouver que c'est une Hot-équivalence. Mais en considérant
\[
u : a \longrightarrow a' \qquad v : b \longrightarrow b'
\]
comme des \textit{chemins} dans $X$, et
\[
u^o : a' \longleftarrow a
\]
comme un chemin rétrograde de $a'$ à $a$, la théorie de variance générale (dans IX) nous donne bel et bien un isomorphisme dans (Hot) :
\[
\Ch_\infty(X;a,b) \longrightarrow \Ch_\infty(X;a',b'),
\]
on a envie de dire que c'est $\varphi$, i.e. qu'en l'insérant à la place de $\varphi$, le carré ci-dessus devient commutatif (ce qui prouverait ce qu'on veut). Je pense que si on travaille avec les ind-objets $\CCh_\infty(X)/(a,b)$, $\CCh_\infty(X;a,b)$ etc., pour lesquels on a des accouplements qui s'explicitent agréablement aux crans finis, et donnent naissance aux $\varphi$, cela devrait être facile (quoiqu'un peu fastidieux). Cela ne prouverait pas pour autant la commutativité de $(*)$ dans Hot, mais
% foot: Version 22 décembre 2001
\printedpage{95}
% head: A. Grothendieck : Dérivateurs, Ch. VII, Édition des Universités de Montpellier II et Paris VII
\msp{page 152}
seulement dans $\Hotz$, ce qui n'est peut-être pas suffisant. Mais cela suffit pour prouver que la flèche $\rho$ dans $(*)$ devient \texttt{[un]} isomorphisme dans $\Hotz$, donc elle est \texttt{[un]} isomorphisme.

2$^o$) Il faut montrer que les HOT-fibrations
\[
\Ch_\infty(X) \overset{\sigma}{\underset{\beta}{\rightrightarrows}} X
\]
sont triviales, ou ce qui revient au même, qu'elles sont à fibres asphériques. Il suffit de le voir pour $\sigma$, donc que
\[
\Ch_\infty(X;a,?)
\]
est asphérique. Mais cette catégorie est la limite inductive des $\Ch_n(X;a,?)$, i.e. des
\[
\underline{\mathrm{Hom}}^{\text{source fixée}}(\tau^n, X).
\]
Mais comme l'intervalle $I_{\tau^n}$ se rétracte par déformation sur son origine, il est immédiat que cette catégorie \underline{Hom} est même contractile.

3$^o$) Considérons
\[
\begin{tikzcd}[column sep=small, row sep=large]
 & \Xb = C(f) \arrow[dl] \arrow[dr, "g"] \arrow[ddrr, bend left, "\bar f"] & & \\
X' \arrow[ur, dotted, bend left, "\alpha"] \arrow[dr, "f"'] & \underline{\text{cart.}} & \Ch_\infty(X) \arrow[dl, "\sigma"'] \arrow[dr, "\beta"] & \\
 & X & & X,
\end{tikzcd}
\]
\msp{page 153}
je voudrais prouver que $\bar f$ est une \textbf{D}-fibration propre et lisse (\footnote[38]{Ça résulte formellement de ce qui a été déjà vu, cf. pages 170-171.}).

$\Xb$ est la catégorie des couples
\[
(x',c) \ \in \ \mathrm{Ob}\,X' \times \mathrm{Ch}_\infty(X)
\]
tels que
\[
\sigma(c) = f(x')
\]
et
\[
\bar f(x',c) = \beta(c), \qquad \text{but de } c.
\]
La fibre de $\Xb$, $\bar f$ en $x_0 \in X$ est donc la catégorie des couples $(x',c) \in \mathrm{Ob}\,X' \times \mathrm{Ch}_\infty(X)$, tels que
\[
\sigma(c) = f(x'), \qquad \beta(c) = x_0.
\]
% foot: Version 22 décembre 2001

\printedpage{96}
% head: A. Grothendieck : Dérivateurs, Ch. VII, Édition des Universités de Montpellier II et Paris VII
On peut regarder encore $\Xb$ comme limite inductive des sous-catégories $\Xb_{n,m}$, correspondante aux couples $(x',c)$ avec $c \in \mathrm{Ch}_{n,m}(X)$. On trouve que pour $m$ pair (direction terminale de $\tau_{n,m}$ rétrograde) les $\Xb_{n,m}$ sont Cat-fibrées sur $X$, pour $m$ impair (direction terminale de $\tau_{n,m}$ directe) elles sont Cat-cofibrées. Donc on conclut encore comme tantôt, $\Xb$ étant (au choix) limite inductive filtrante de sous-catégories lisses sur $X$, ou propres sur $X$. % [sic] correspondante -> correspondantes, agreement with « sous-catégories »

\msp{page 154}
Il faut prouver enfin que $\Xb \xrightarrow{\bar f} X$ est une \textbf{D}-fibration (\footnote[39]{Cela résulte déjà du fait que c'est propre et lisse!}), ce qui revient p. ex. à dire que pour une flèche
\[
a \xrightarrow{v} b
\]
dans $X$, la flèche induite
\[
\Xb/a \xrightarrow{\bar v} \Xb/b
\]
est une Hot-équivalence. On sait déjà comme tantôt que l'on a un diagramme à flèches verticales des Hot-équivalences
\[
\begin{tikzcd}[row sep=large]
\Xb/a \arrow[r] & \Xb/b \\
\Xb_a \arrow[u, hook, "\wr"] \arrow[r, dotted] & \Xb_b, \arrow[u, hook, "\wr"]
\end{tikzcd}
\]
et tout revient à montrer qu'une certaine flèche entre $\Xb_a$ et $\Xb_b$ dans Hot est \texttt{[un]} isomorphisme. J'ai envie de définir directement une telle flèche dans Hot, voire dans Cat. Or pour les $m$ impairs, les $\Xb_{n,m}$ sont Cat-cofibrées sur $X$. Pour $\Xb_{-m,m}$ ($m$ impair) on trouve donc
\[
(\Xb_{-m,m})_a \xrightarrow{v_*} (\Xb_{-m,m})_b,
\]
\msp{page 155}
l'ennui est qu'elles ne sont pas compatibles avec les morphismes de transition pour $m$ variable. Elles le sont cependant \textit{à homotopie près} sûrement, donc on trouvera au moins une flèche d'ind-objets dans Cat\,hot
\[
\Xbb_a \xrightarrow{v_*} \Xbb_b,
\]
ce qui par passage à $\Hotz$ donne quand-même une flèche
\[
\Xb_a \xrightarrow{v_*} \Xb_b.
\]
En travaillant avec les $m$ impairs, on trouve de même des foncteurs en sens inverses % [sic] m impairs -> m pairs, the Cat-fibrées (giving v^*) are stated above for m pair
\[
(\Xb_{-m,m})_a \xleftarrow{v^*} (\Xb_{-m,m})_b,
\]
% foot: Version 22 décembre 2001
\printedpage{97}
% head: A. Grothendieck : Dérivateurs, Ch. VII, Édition des Universités de Montpellier II et Paris VII
qui ne se recollent pas pour $m$ variable, mais qui se recollent à homotopie près, pour donner une flèche dans Ind\,Cat\,hot
\[
\Xbb_a \xleftarrow{v^*} \Xbb_b,
\]
d'où aussi
\[
\Xb_a \xleftarrow{v^*} \Xb_b \qquad \text{dans } \Hotz.
\]
D'autre part, on doit pouvoir montrer soritalement,
\msp{page 156}
cran par cran, que $v_*v^*$ et $v^*v_*$ dans Ind\,Cat\,hot sont égales à l'identité, donc $v_*$ et $v^*$ sont des isomorphismes inverses l'un de l'autre. On doit vérifier d'autre part soritalement que
\[
\begin{tikzcd}[row sep=large]
\Xbb/a \arrow[r, "\bar v"] & \Xbb/b \\
\Xbb_a \arrow[u] \arrow[r, "v_*"] & \Xbb_b, \arrow[u]
\end{tikzcd}
\]
dans Ind\,Cat\,hot est commutatif. Passant au diagramme correspondant dans $\Hotz$, en prenant les limites inductives, on trouve un carré commutatif dans $\Hotz$
\[
\begin{tikzcd}[row sep=large]
\Xb/a \arrow[r, "\bar v"] & \Xb/b \\
\Xb_a \arrow[u, "\wr"'] \arrow[r, "v_*", "\sim"'] & \Xb_b, \arrow[u, "\wr"']
\end{tikzcd}
\]
d'où résulte que $\bar v$ est \texttt{[un]} isomorphisme dans $\Hotz$, donc aussi dans Hot.

Pour vérifier que la HOT-fibration $\bar f$ est triviale si et seulement si $f$ est une HOT-équivalence, il faudrait avoir développé dans Cat le formalisme de la suite exacte des lacets - on est tout près visiblement !
% foot: Version 22 décembre 2001

\printedpage{98}
% head: A. Grothendieck : Dérivateurs, Ch. VII, Édition des Universités de Montpellier II et Paris VII
\msp{page 157}

\begin{flushright}
\footnotesize 13.11.90
\end{flushright}

\section*{13\quad Où on prouve enfin que $\Ch(X;-,-)$ est aussi bon (voire meilleur) que $\Ch_\infty(X;-,-)$.}

Je crois finalement que les catégories $\Ch(X)$, $\Ch(X;a,b)$ (sans indice $\infty$!) sont `bonnes', que les flèches canoniques dans Hot
\[
\left\{
\begin{aligned}
\Ch_\infty(X) \quad &\xrightarrow{\sim} \quad \Ch(X)\\
\Ch_\infty(X;a,b) \quad &\xrightarrow{\sim} \quad \Ch(X;a,b)
\end{aligned}
\right.
\]
sont des isomorphismes. (La première est même \texttt{[un]} isomorphisme dans $\mathrm{HOT}(X \times X)$, ce qui implique la deuxième relation, en passant aux fibres \dots) Je crois que je suis à présent (après XI) en mesure de le prouver. Cela résultera des deux résultats-clef suivants :

\begin{enumerate}
\item[1)] Le foncteur canonique
\[
X \longrightarrow \Ch(X)
\]
(associant à tout $x \in \mathrm{Ob}\,X$ le chemin vide de source et but $X$) est une Hot-équivalence, et : % [sic] but $X$ -> but $x$, the object is $x \in \mathrm{Ob}\,X$
\item[2)] Le foncteur source-but
\[
\Ch(X) \longrightarrow X \times X
\]
est \textit{lisse} et \textit{propre} (donc une `Hot-fibration parfaite').
\end{enumerate}

Cela doit donner ce qu'on cherche. Comme la $X \times X$-flèche $\Ch_\infty(X) \longrightarrow \Ch(X)$ (et les flèches déduites $\Ch_\infty(X;a,b) \longrightarrow \Ch(X;a,b)$) a été définie par l'intermédiaire de la catégorie
\[
\mathcal{C}_{\mathbf{N}^o} = \mathbf{N}^o \times_{\Ty} \mathcal{C}
\qquad
\begin{aligned}
&\text{(où $\mathcal{C} = \Ch(X)$ \texttt{[est]} considérée comme catégorie fibrée sur $\Ty$,}\\
&\text{et $\mathbf{N}^o \longrightarrow \Ty$ est le système projectif des $\underbrace{(\Delta_+\Delta_-)^{2n}}_{t_n} \dots$)}
\end{aligned}
\]
\msp{page 158}
à l'aide des deux flèches canoniques
\[
\begin{tikzcd}[row sep=huge, column sep=large]
\mathcal{C}_{\mathbf{N}^o} \arrow[r] \arrow[d] & \Ch_\infty(X) \\
\Ch(X), &
\end{tikzcd}
\]
% foot: Version 22 décembre 2001

% ---- from Der-ChVII-456_p101-105.tex
\printedpage{99}
% head: A. Grothendieck : Dérivateurs, Ch. VII, Édition des Universités de Montpellier II et Paris VII
on sera obligé de prouver aussi les résultats qui correspondent à 1) 2) pour $\cC_{\NN}$, i.e.

\begin{itemize}[label={}]
\item[1$'$)]
\[
\begin{tikzcd}
X \ar[r] \ar[d,"\wr"'] & \cC_{\NN} \\
\cC_{t_0} = \cC_0 &
\end{tikzcd}
\]
(l'inclusion canonique de $\cC_0$, fibre de $\cC_{\NN}$ (fibrée sur $\mathbf{N}^o$) en $0 \in \mathbf{N}^o$) est une Hot-équivalence, et :

\item[2$'$)] $\cC_{\NN} \longrightarrow X \times X$ est propre et lisse.
\end{itemize}

Moyennant ces quatre résultats techniques, que je vais prouver tantôt (et $1'$ et $2'$ ne font aucune surprise pour moi), on gagne. Considérons en effet le diagramme
\[
\begin{tikzcd}
& X \ar[dl,"\alpha_1"'] \ar[d,"\alpha_0"] \ar[dr,"\alpha_2"] & \\
\Ch(X) \ar[dr,"p_1"'] & \cC_{\NN} \ar[l,"\pi_1"'] \ar[r,"\pi_2"] \ar[d,"p_0"] & \Ch_{\infty}(X) \ar[dl,"p_2"] \\
& X \times X. &
\end{tikzcd}
\]

Par 1) et $1'$) (concernant $\alpha_1$ et $\alpha_0$) et le résultat
\msp{page 159}
similaire déjà connu pour $\alpha_2$, ces trois foncteurs sont des Hot-équivalences. Donc il en est de même de $\pi_1$, $\pi_2$. Comme $p_1$, $p_0$, $p_2$ sont des Hot-fibrations (étant propres et lisses), il s'ensuit que $\pi_1$ et $\pi_2$ induisent des Hot-équivalences sur les fibres
\[
\Ch(X;a,b) \quad \xleftarrow{\ \pi_1(a,b)\ } \quad \cC_{\NN}(a,b) \quad \xrightarrow{\ \pi_2(a,b)\ } \quad \Ch_{\infty}(X;a,b).
\]
Je l'avais vérifié (sans mal) déjà pour $\pi_2$ et les $\pi_2(a,b)$, ce qui avait permis, en inversant $\pi_2$ dans Hot, de définir \textit{dans} Hot
\[
\left\{
\begin{aligned}
\Ch_{\infty}(X) &\quad\longrightarrow\quad \Ch(X) \\
\Ch_{\infty}(X;a,b) &\quad\longrightarrow\quad \Ch(X;a,b),
\end{aligned}
\right.
\]
et dire que ce sont des isomorphismes signifie justement que $\pi_1$ et les $\pi_1(a,b)$ sont des isomorphismes.

Je ne crois pas que $\cC_{\NN}$ ait en lui-même d'autre interêt, que de permettre de comparer % [sic] interêt -> intérêt, missing acute accent
les $X \times X$-catégories $\Ch_{\infty}(X)$ et $\Ch(X)$, à l'aide du diagramme précédent. Par contre, $\Ch_{\infty}(X)$ et $\Ch(X)$ m'aparaissent comme les deux ingrédients essentiels du formalisme des % [sic] m'aparaissent -> m'apparaissent, spelling
chemins dans Cat.
% foot: Version 22 décembre 2001

\printedpage{100}
% head: A. Grothendieck : Dérivateurs, Ch. VII, Édition des Universités de Montpellier II et Paris VII
\msp{page 160}
Je vais prouver 1, 2, $1'$, $2'$.

1) Considérons le diagramme
\[
\begin{tikzcd}
X = (X_{\Ty})_e \ar[rr,hook,"i"] \ar[dr,dotted] & & X_{\Ty} = X \times \Ty \ar[rr,"j"] \ar[dr] & & \Ch(X) \ar[dl] \\
& e \ar[rr,dotted] & & \Ty\ , &
\end{tikzcd}
\]
où $X_{\Ty}$ et $\Ch(X)$ sont considérées comme catégories Cat-fibrées sur $\Ty$, et le foncteur $j$ est déduit des foncteurs canoniques
\[
X \longrightarrow \Ch(X)_{\tau} \overset{\text{déf}}{=} \Ch_{\tau}(X),
\]
associant à $x \in \Ob X$ le chemin $p_{\tau}^{*}(x)$, où $p_{\tau} : \tau \longrightarrow e$ est la projection de $\tau$ vers l'élément unité (et élément final) $e = 1_{\Ty}$ de $\Ty$, et $\Ch_{e}(X)$ est identifiée à $X$. Il faut vérifier le

\textbf{Lemme.} \textit{Pour tout $\tau \in \Ty$, le foncteur précédent est une} Hot-\textit{équivalence (et même un homotopisme).}

Mais on sait que pour deux catégories $\cX$, $\cC$ lisses sur une autre, $T$, tout $T$-morphisme de $\cX$ dans $\cC$ qui induit des Hot-équivalences sur les fibres, est une Hot-équivalence. Ainsi $j$ est
\msp{page 161}
est une Hot-équivalence. D'autre part, $\Ty$, ayant un objet final, est asphérique, donc % [sic] est ... est -> est, word repeated across the manuscript page break
$X \times \Ty \xrightarrow{\mathrm{pr}_1} X$ est une Hot-équivalence, d'où résulte aussitôt que $i$ est une Hot-équivalence. Donc de même $ji$, or $ji = \alpha_1$, q.e.d.

2) Prouvons p. ex. la lissité de $\Ch(X) \longrightarrow X \times X$, la propriété s'en déduira en appliquant ce résultat à $X^{\mathrm{opp}}$. Soit donc
\[
(a',b') \ \longrightarrow\ (a,b), \qquad \text{i.e. } a' \xrightarrow{u} a,\ \ b' \xrightarrow{v} b,
\]
une flèche dans $X \times X$, et $c$ un objet de $\Ch(X)$ au dessus de $(a,b)$ :
\[
\begin{tikzcd}[column sep=large]
a' \ar[r,dotted,"c'","\scriptstyle\rightarrow"'] \ar[d,"u"'] & b' \ar[d,"v"] \\
a \ar[r,"c","\scriptstyle\rightarrow"'] & b.
\end{tikzcd}
\]

On regarde la catégorie $\cC(c;u,v)$ des couples $(c',w)$ qui relèvent $(c,(u,v))$, où donc $c'$ est un chemin (de type non précisé) de $a'$ vers $b'$, et
\[
w : c' \longrightarrow c
\]
une flèche dans $\Ch(X)$ qui sur les objets source induise $u$, et sur les objets buts induise $v$ (ainsi $c'$ est un chemin de $a'$ vers $b'$). Il faut
% foot: Version 22 décembre 2001
\printedpage{101}
% head: A. Grothendieck : Dérivateurs, Ch. VII, Édition des Universités de Montpellier II et Paris VII
\msp{page 162}
prouver que cette catégorie $\cC(c;u,v)$ est asphérique. Pour cela, on regarde le chemin composé
\[
\tilde{c} = \mathrm{ch}(u)\ c\ \mathrm{ch}(v)^{o},
\]
où $\mathrm{ch}(u)$ est le chemin (direct) de $a'$ vers $a$ défini par $u : a' \longrightarrow a$, et $\mathrm{ch}(v)^{o}$ désigne donc le chemin (rétrograde) de $b$ à $b'$ défini par $v : b' \longrightarrow b$. Si $c$ est de type $\tau$, $\tilde{c}$ est donc de type $\Delta_{+}\tau\Delta_{-}$. On va comparer $\cC(c;u,v)$ et la catégorie $\Ch(X)/\tilde{c}$, à l'aide de deux foncteurs
\[
\begin{tikzcd}
\cC(c;u,v) \ar[r,shift left,"\varphi"] & \Ch(X;a',b')/\tilde{c}\ , \ar[l,shift left,"\psi"]
\end{tikzcd}
\]
dont on prouvera qu'ils définissent des flèches inverses l'une de l'autre dans Hot. Comme $\Ch(X)/\tilde{c}$ est contractile (ayant un objet final), donc asphérique, il en est de même de $\cC(c;u,v)$, et on gagnera.

On définit $\varphi$ par $\boxed{\varphi(c',w') = (\tilde{c}',\tilde{w}'),}$ où ($c'$ étant du type $\tau'$) $\tilde{c}'$ est le chemin de type $\Delta_{+}\tau'\Delta_{-}$ déduit de $c'$ comme image inverse par la dégénérescence $\Delta_{+}\tau'\Delta_{-} \longrightarrow \tau'$, i.e. on prolonge $c'$ d'un `cran $\Delta_{+}$' $a' \xrightarrow{\mathrm{id}} a'$ à gauche et d'un `cran $\Delta_{-}$' $b' \xleftarrow{\mathrm{id}} b'$ à droite, cf. figure.
\[
\overbrace{\underbrace{\begin{tikzcd}[ampersand replacement=\&]
a' \ar[r,"\mathrm{id}"] \ar[d,"\mathrm{id}"'] \ar[dr,phantom,"\ast"] \& a' \ar[rr,dotted,no head,"c'"] \ar[d,"u"'] \& {} \ar[d,dotted,"w'"] \& b' \ar[d,"v"'] \ar[dr,phantom,"\ast"] \& b' \ar[l,"\mathrm{id}"'] \ar[d,"\mathrm{id}"] \\
a' \ar[r,"u"] \& a' \ar[rr,no head,"c'"'] \& {} \& b' \& b' \ar[l,"v"'] % [sic] a' ... b' with label c' in the bottom row -> a ... b with label c, the bottom row of the figure is the path c~ = ch(u) c ch(v)^o
\end{tikzcd}}_{\tilde{c}}}^{\tilde{c}'}
\]

\msp{page 163}
On peut dire que
\[
\tilde{c}' = a'_{\Delta_{+}} \circ c' \circ b'_{\Delta_{-}},
\]
où de façon générale, si $x \in X$ et $\tau \in \mathrm{Ty}$, $x_{\tau}$ désigne le `chemin constant' de valeur $x$, type $\tau$. On a d'autre part des morphismes évidents
\[
a'_{\Delta_{+}} \longrightarrow \mathrm{ch}(u), \qquad b'_{\Delta_{-}} \longrightarrow \mathrm{ch}(v)^{\mathrm{opp}},
\]
comme montré [?] par les flèches verticales $(\mathrm{id}_{a'},u)$ et $(v,\mathrm{id}_{b'})$ dans les deux carrés marqués $\ast$ du diagramme, flèches qui sont l'une dans $\Ch_{\Delta_{+}}(X)$, l'autre dans $\Ch_{\Delta_{-}}(X)$. Ceci dit, on pose
\[
\tilde{w}' : \tilde{c}' \longrightarrow \tilde{c}, \qquad \tilde{w} = \lambda(u) \ast w' \ast \mu(v), % [sic] \tilde{w} -> \tilde{w}', the morphism just named is \tilde{w}'
\]
où
\[
\left\{
\begin{aligned}
\lambda(u) :\ \ a'_{\Delta_{+}} &\quad\longrightarrow\quad \mathrm{ch}(u) \\
\mu(v) :\ \ b'_{\Delta_{-}} &\quad\longrightarrow\quad \mathrm{ch}(v)^{\mathrm{opp}}
\end{aligned}
\right.
\]
% foot: Version 22 décembre 2001
\printedpage{102}
% head: A. Grothendieck : Dérivateurs, Ch. VII, Édition des Universités de Montpellier II et Paris VII
(cf. figure). Le couple
\[
(\tilde{c}',\tilde{w}') \in \Ob \Ch(X)/\tilde{c}
\]
ainsi défini dépend fonctoriellement de $(c',w')$ variable dans $\cC(c;u,v)$ : Si $(c'_1,w'_1)$ est un deuxième objet de cette catégorie, une flèche du premier dans le second est une flèche
\[
c' \xrightarrow{\alpha} c'_1
\]
dans $\Ch(X;a',b')$, telle qu'on ait commutativité dans
\[
\begin{tikzcd}
c' \ar[rr,"\alpha"] \ar[dr,"w'"'] & & c'_1 \ar[dl,"w'_1"] \\
& c, &
\end{tikzcd}
\qquad\qquad w' = w'_1\alpha.
\]

Une telle flèche $\alpha$ définit
\[
\tilde{\alpha} = a'_{\Delta_{+}} \ast \alpha \ast b'_{\Delta_{-}} : \tilde{c}' \longrightarrow \tilde{c}'_1
\]
dans $\Ch(X;a',b')$, et on vérifie
\msp{page 164}
aussitôt que l'on a commutativité dans
\[
\begin{tikzcd}
\tilde{c}' \ar[rr,"\tilde{\alpha}"] \ar[dr,"\tilde{w}'"'] & & \tilde{c}'_1 \ar[dl,"\tilde{w}'_1"] \\
& \tilde{c}, &
\end{tikzcd}
\qquad\qquad \text{i.e. } \tilde{w}' = \tilde{w}'_1\tilde{\alpha},
\]
i.e. que $\tilde{\alpha}$ est bien une flèche dans $\Ch(X;a',b')/\tilde{c}$. De même, que $\alpha \longmapsto \tilde{\alpha}$ définit une dépendance fonctorielle, i.e. $(\beta\alpha)\tilde{} = \tilde{\beta}\tilde{\alpha}$, est évident sur la définition de $\tilde{\alpha}$. Donc, on a bien défini % Reading resolved against the supplied PDF; see corrections.csv.
\[
\cC(c;u,v) \ \xrightarrow{\ \varphi\ }\ \Ch(X;a',b')/\tilde{c},
\]
définissons maintenant
\[
\cC(c;u,v) \ \xleftarrow{\ \psi\ }\ \Ch(X;a',b')/\tilde{c}.
\]
Soit $(c'',w'')$ un objet de $\Ch(X;a',b')/\tilde{c}$,
\[
\overbrace{\underbrace{\begin{tikzcd}[ampersand replacement=\&]
a'_0 = a' \ar[r] \ar[d,"\mathrm{id}"'] \& a'_1 \ar[rr,no head,"\rightarrow"] \ar[d,"w''_1"'] \& {} \ar[d,dotted,"w''"] \& a'_n \ar[d,"w''_n"'] \& b' \ar[l] \ar[d,equal] \\
a' \ar[r,"u"] \& a \ar[rr,no head] \& {} \& b \& b', \ar[l,"v"']
\end{tikzcd}}_{c}}^{c''}
\]
% foot: Version 22 décembre 2001

\printedpage{103}
% head: A. Grothendieck : Dérivateurs, Ch. VII, Édition des Universités de Montpellier II et Paris VII
donc
\[
w'' : c'' \longrightarrow \tilde{c}, \qquad\quad \text{flèche dans } \Ch(X;a',b'),
\]
notons qu'on a une flèche de rétraction évidente
\[
\tilde{c} \xrightarrow{\rho_c} c \qquad\quad \text{associée à } \Delta_{+}\tau\Delta_{-} \longrightarrow \tau \text{ (rétraction),}
\]
induisant $u$ et $v$ sur la source et le but :
\[
\begin{tikzcd}
a' \ar[r,"u"] \ar[d,"u"] & a \ar[rrr,no head,"c"] \ar[d,"\mathrm{id}"] & {} \ar[d,"\mathrm{id}"] & {} \ar[d,"\mathrm{id}"] & b \ar[d,"\mathrm{id}"] & b' \ar[l,"v"'] \ar[d,"v"] & \Big\}\ \tilde{c} & \in \Ch_{\Delta_{+}\tau\Delta_{-}}(a',b') \\
a \ar[r,"\mathrm{id}"] & a \ar[rrr,no head,"c"] & {} & {} & b & b' \ar[l,"v"'] & \Big\}\ \delta^{*}_{\tau}(c) & \in \Ch_{\Delta_{+}\tau\Delta_{-}}(a,b) \\
& a \ar[rrr,no head,"c"] & & & b & & \Big\}\ c & \in \Ch_{\tau}(a,b).
\end{tikzcd}
\]

\msp{page 165}
Donc on définit un objet
\[
(c',w') = \psi(c'',w'')\ \in \Ob \cC(c;u,v)
\]
par
\[
c' = c'', \qquad\quad w' = \rho_c w'',
\]
et c'est fonctoriel en $(c'',w'')$.

Le composé $\varphi\psi$ étant un endomorphisme de la catégorie contractile $\Ch(X;a',b')/\tilde{c}$, est homotope à l'identité. Il faut vérifier qu'il en est de même de $\psi\varphi$, endomorphisme de $\cC(c;u,v)$. Ce foncteur associe à $(c',w')$, $(\tilde{c}',\overline{w'})$, où $\tilde{c}'$ est défini comme plus haut, $\tilde{c}' = a'_{\Delta_{+}} \circ c' \circ b'_{\Delta_{-}}$, et où $\overline{w'}$ est le composé
\[
\underbrace{\tilde{c}'}_{= \delta^{\prime *}_{\tau'}(c')}\ \ \underbrace{\xrightarrow{\ \rho_{c'}\ }}_{\substack{\text{flèche}\\ \text{canonique}}}\ \ c' \xrightarrow{w'} c.
\]

On a donc
\[
(\psi\varphi)(c',w') = (\tilde{c}',\overline{w'}) \xrightarrow{\rho_{c'}} (c',w'),
\]
c'est bien un morphisme dans $\cC(c;u,v)$, i.e. on a commutativité dans
\[
\begin{tikzcd}
\tilde{c}' \ar[rr,"\rho_{c'}"] \ar[dr,"\overline{w'}"'] & & c' \ar[dl,"w'"] \\
& c &
\end{tikzcd}
\qquad\qquad \text{(flèche de } \Ch(X;a',b')\text{).}
\]
% foot: Version 22 décembre 2001

% ---- from Der-ChVII-456_p106-109.tex
\printedpage{104}
% head: A. Grothendieck : Dérivateurs, Ch. VII, Édition des Universités de Montpellier II et Paris VII
\msp{page 166}

Ainsi, le foncteur $\psi\varphi$ est bien homotope au foncteur identique. On a gagné !

Passons à $1'$, $2'$) pour $\CN$. L'assertion $1'$) se démontre exactement de la même manière que 1) pour \underline{Ch}$(X)$, en regardant $\CN$ comme fibrée sur $\mathbf{N}^o$ (qui a un objet final 0), et en utilisant le diagramme
\[
\begin{tikzcd}
X = (X_{\mathbf{N}^o})_e \arrow[rr, "i'"] & & X_{\mathbf{N}^o} = X \times \mathbf{N}^o \arrow[rr, "j'"] \arrow[dr] & & \mathcal{C}_{\mathbf{N}^o} \arrow[dl] \\
 & & & \mathbf{N}^o \ , &
\end{tikzcd}
\]
où $j'$ est (cartésien et) une Hot-équivalence sur chaque fibre d'après le lemme, donc une Hot-équivalence globalement, et $i'$ est une Hot-équivalence car $\mathbf{N}^o$ est asphérique (ayant un objet final). Donc $j'i'$ est asphérique, q.e.d.

Il reste à prouver que
\[ \CN \longrightarrow X \times X \]
est propre et lisse. Ici encore,
\msp{page 167}
c'est exactement le même argument que tantôt, je présume. Il y a cependant une complication technique. Nous avions pris, pour définir
\[ \underline{t} : \mathbf{N}^o \longrightarrow \underline{\mathrm{Ty}}, \qquad \underline{t} = (t_n)_n, \]
le choix
\[ t_n = (\Delta_+\Delta_-)^{2n}, \]
ou si on préfère
\[ t_n = (\Delta_+\Delta_-)^n, \]
en tous cas les $t_n$ sont à direction initiale $+$, terminale $-$, et ce sont des types réduits. Si on prend $\Delta_+t_n\Delta_-$, on trouve un type non réduit, qui n'est pas parmi les $t_n$. Il ne servirait à rien de changer la définition de $t_n$ en $t_n = (\Delta_-\Delta_+)^{\text{quelque chose}}$, car cette fois $\Delta_+t_n\Delta_-$ sera bien réduit, mais n'aura pas les bons signes initial et terminal. Donc pour définir $\tilde{c}$ en termes de $c$, comme objet de $\CN$, il faudra, si $c$ est de type $\tau = t_n$, prendre $\tilde{c}$ de type $\Delta_+\Delta_-\tau\Delta_+\Delta_- = t_{n+1}$ (en prenant $t_{n+1} = (\Delta_+\Delta_-)^{2n}$), et de même pour définir $\tilde{c}'$ en termes de $c'$. Les définitions pertinentes du passage de $(c',w')$ à $(\tilde{c}',\tilde{w}')$ sont décrites dans la figure : % [sic] t_{n+1} = (\Delta_+\Delta_-)^{2n} -> t_{n+1} = (\Delta_+\Delta_-)^{2n+2}, exponent inconsistent with t_n = (\Delta_+\Delta_-)^{2n}

\msp{page 168}
% foot: Version 22 décembre 2001

\printedpage{105}
% head: A. Grothendieck : Dérivateurs, Ch. VII, Édition des Universités de Montpellier II et Paris VII
\[
\tilde{w}'\Big\downarrow
\left\{
\overbrace{\underbrace{
\begin{tikzcd}[ampersand replacement=\&]
a' \arrow[r, "\mathrm{id}"] \arrow[d, "\mathrm{id}"'] \& a' \arrow[d, "u"'] \& a' \arrow[l, "\mathrm{id}"'] \arrow[d, "u"'] \arrow[rrrr, dotted, no head, "c'", "\longrightarrow"'] \& {} \arrow[d, "w'"] \& {} \arrow[d, "w'"] \& {} \arrow[d, "w'"] \& b' \arrow[r, "\mathrm{id}"] \arrow[d, "v"'] \& b' \arrow[d, "v"] \& b' \arrow[l, "\mathrm{id}"'] \arrow[d, "\mathrm{id}"] \\
a' \arrow[r, "u"] \& a \& a \arrow[l, "\mathrm{id}"'] \arrow[rrrr, no head, "c", "\longrightarrow"'] \& {} \& {} \& {} \& b \arrow[r, "\mathrm{id}"] \& b \& b'. \arrow[l, "v"']
\end{tikzcd}
}_{\tilde{c}}}^{\tilde{c}'}
\right\}
\tilde{w}'
\]

Ainsi
\[
\dcfit{\left\{
\begin{array}{lcl ll}
\tilde{c} & = & \xi_u \circ c \circ \eta_v & \text{où} &
\left\{
\begin{array}{lccccc}
\xi_u : & a' & \xrightarrow{u} & a & \xleftarrow{\mathrm{id}} & a \\[1ex]
\eta_v : & b & \xrightarrow{\mathrm{id}} & b & \xleftarrow{v} & b'
\end{array}
\right\}
\text{ type } \Delta_+\Delta_- \\[4ex]
\tilde{c}' & = & a'_{\Delta_+\Delta_-} \circ c' \circ b'_{\Delta_+\Delta_-} & & \\[4ex]
\tilde{w}' & = & \lambda(u) * w * \mu(v) & \text{avec} & % [sic] w -> w', the homomorphism of the pair (c',w')
\left\{
\begin{array}{lccc}
\lambda(u) : & a'_{\Delta_+\Delta_-} & \longrightarrow & \xi_u \\[1ex]
\lambda(v) : & b'_{\Delta_+\Delta_-} & \longrightarrow & \eta_v % [sic] \lambda(v) -> \mu(v), the formula above has \mu(v)
\end{array}
\right\}
\text{ comme dans \texttt{[la]} figure.}
\end{array}
\right.}
\]
On trouve ainsi, mutatis mutandis,
\[ \CN(c;u,v) \xrightarrow{\varphi} \CN(a',b')/\tilde{c}, \qquad \varphi(c',w') = (\tilde{c}',\tilde{w}'), \]
et on définit le foncteur
\[ \CN(c;u,v) \xleftarrow{\psi} \CN(a',b')/\tilde{c} \]
ainsi : à $(c'',w'')$ dans $\CN(a',b')/\tilde{c}$, donc un chemin $c''$ de type $t_n = (\Delta_+\Delta_-)^{2n}$ de $a'$ à $b'$ et un homomorphisme $w'' : c'' \longrightarrow \tilde{c}$ à extrémités fixes, on associe $(c',w')$ dans $\CN(c;u,v)$, avec
\[ c' = c'', \qquad w' = \rho_c w'', \]
où
\[ \rho_c : \tilde{c} \longrightarrow c \]
est défini de façon similaire que tantôt, par
\msp{page 169}
\[
\begin{tikzcd}
a' \arrow[r, "u"] \arrow[d, "u"'] & a \arrow[d, "\mathrm{id}"] & a \arrow[l, "\mathrm{id}"'] \arrow[d, "\mathrm{id}"] \arrow[rr, zigzag] & {} \arrow[d, "\mathrm{id}"'] & b \arrow[r, "\mathrm{id}"] \arrow[d, "\mathrm{id}"] & b \arrow[d, "\mathrm{id}"] & b' \arrow[l, "v"'] \arrow[d, "v"] & \tilde{c},\ \text{ type } t_{n+1}. \\
a \arrow[r, "\mathrm{id}"] & a & a \arrow[l, "\mathrm{id}"'] \arrow[rr, zigzag, "c"] & {} & b \arrow[r, "\mathrm{id}"] & b & b' \arrow[l, "\mathrm{id}"'] & \varphi^*_{n,n+1}(c).
\end{tikzcd}
\]

On doit seulement vérifier que l'endofoncteur $\psi\varphi$ de $\CN(c;u,v)$ est homotope à l'identité. Mais ce foncteur transforme $(c',w')$ en $(\tilde{c}' = a'_{\Delta_+\Delta_-}c'b'_{\Delta_+\Delta_-}, \overline{w'})$, où $\overline{w'}$ est encore le composé
% foot: Version 22 décembre 2001
\printedpage{106}
% head: A. Grothendieck : Dérivateurs, Ch. VII, Édition des Universités de Montpellier II et Paris VII
\[
\begin{tikzcd}
\tilde{c}' \arrow[r, "\text{can.}"] \arrow[rr, bend right=20, "\overline{w'}"'] & c' \arrow[r, "w'"] & c.
\end{tikzcd}
\]

On a donc
\[ \psi\varphi \longrightarrow \mathrm{id}, \]
$\psi\varphi$ est homotope au foncteur identique.

Cela prouve la lissité de $\CN \longrightarrow X \times X$. Pour la propreté, il faut prouver la lissité de le cas de $X^o$, et du système projectif $\underline{t}^{\mathrm{opp}} \simeq (t_n^{\mathrm{opp}})_n$. Mais on a $\underline{t}^{\mathrm{opp}} = \underline{t}$, et il suffit d'appliquer ce qu'on a déjà prouvé au cas de $X^o$, pour voir que $\CN \longrightarrow X \times X$ est propre. % [sic] de le cas -> dans le cas, grammatical

Il ne reste plus que le lemme à prouver. Ça doit être essentiellement
\msp{page 170}
trivial, et je renonce à rédiger une démonstration.

Je vais maintenant montrer `formellement' que la factorisation `canonique' d'une flèche de Cat $f : X' \longrightarrow X$ en
\[ X' \xrightarrow{i_f} \overline{X'} \xrightarrow{\bar{f}} X, \qquad i_f \text{ Hot-équivalence,} \]
avec (on l'espérait) $\bar{f}$ propre et lisse, à partir d'une factorisation similaire pour la flèche diagonale $X \xrightarrow{\ \ \mathrm{diag}_X\ \ } X \times X$ en
\[
\begin{tikzcd}
X \arrow[d, "i_X\ \ \text{\footnotesize Hot-équivalence}"] \\
\mathcal{C}_X\ \ \text{`chemins'} \arrow[d, "p_X"] \\
X \times X
\end{tikzcd}
\]
est chose formelle. (Cf. page 139 pour les notations et \texttt{[la]} problématique.) On regarde le diagramme
\[
\begin{tikzcd}
 & \mathcal{C}_X \arrow[d, "p_X"] \arrow[dl, "\beta_X"'] \arrow[dr, phantom, "\text{cart.}"] & C(f) \arrow[l, "\mathrm{pr}'_2"'] \arrow[d, "p_{X,f}"] \arrow[dll, bend right=50, "\bar{f}"'] & [-1.5em] \simeq X' \times_X (C_X, \sigma = \mathrm{pr}_1 \circ p_X) \\ % Reading resolved against the supplied PDF; see corrections.csv.
X & X \times X \arrow[l, "\mathrm{pr}_2^X"] \arrow[d, "\mathrm{pr}_1^X"'] \arrow[dr, phantom, "\text{cart.}"] & X' \times X \arrow[l, "f \times \mathrm{id}_X = f'"] \arrow[d, "\mathrm{pr}_1"] & \\
 & X & X' \arrow[l, "f"'] \arrow[uu, bend right=60, dotted, "i_f"'] &
\end{tikzcd}
\]
% foot: Version 22 décembre 2001
\printedpage{107}
% head: A. Grothendieck : Dérivateurs, Ch. VII, Édition des Universités de Montpellier II et Paris VII
(\footnote[40]{NB $C(f)$ est défini comme l'image inverse sur $X'$ du fibré $C_X$ sur $X$ (par $\sigma_X = \mathrm{pr}_1^X \circ p_X$) par $f : X' \longrightarrow X$.}). Comme $p_X$ est propre et lisse, i.e. `parfait', son image inverse sur $X' \times X$, savoir $p_{X,f}$, l'est aussi; d'autre part $\mathrm{pr}_2^X \circ f'$ est la
\msp{page 171}
projection $\mathrm{pr}_2$ de $X' \times X$ sur $X$, donc c'est aussi un morphisme parfait (puisque $X' \longrightarrow e$ est parfait, donc il en est de même du composé
\[ (\mathrm{pr}_2^X \underbrace{f') \circ p_{X,f}} = \underbrace{\mathrm{pr}_2^X(p_X}\, \mathrm{pr}'_2) = \beta_X \mathrm{pr}'_2 \overset{\text{déf}}{=} \bar{f} \ \ ). \]

Je rappelle que la section $i_f$ de $C(f)$ sur $X'$ est définie à partir de la section $i_X$ de $(\mathcal{C}_X,\sigma_X)$ sur $X$. C'est une Hot-équivalence, car $\sigma_X = \mathrm{pr}_1^X p_X$ est parfait à fibres asphériques, donc universellement une Hot-équivalence, donc $\mathrm{pr}_1 p_{X,f} : C(f) \longrightarrow X'$ est aussi une Hot-équivalence, et toute section en est donc une aussi.
% foot: Version 22 décembre 2001
\end{document}
